\chapter{Advanced Topics}\label{svn.advanced}

If you\textquotesingle ve been reading this book chapter by chapter,
from start to finish, you should by now have acquired enough knowledge
to use the Subversion client to perform the most common version control
operations. You understand how to check out a working copy from a
Subversion repository. You are comfortable with submitting and receiving
changes using the \texttt{svn\ commit} and \texttt{svn\ update}
operations. You\textquotesingle ve probably even developed a reflex that
causes you to run the \texttt{svn\ status} command almost unconsciously.
For all intents and purposes, you are ready to use Subversion in a
typical environment.

But the Subversion feature set doesn\textquotesingle t stop at ``common
version control operations.'' It has other bits of functionality besides
just communicating file and directory changes to and from a central
repository.

This chapter highlights some of Subversion\textquotesingle s features
that, while important, may not be part of the typical
user\textquotesingle s daily routine. It assumes that you are familiar
with Subversion\textquotesingle s basic file and directory versioning
capabilities. If you aren\textquotesingle t, you\textquotesingle ll want
to first read \hyperref[svn.basic]{Fundamental Concepts} and
\hyperref[svn.tour]{Basic Usage}. Once you\textquotesingle ve mastered
those basics and consumed this chapter, you\textquotesingle ll be a
Subversion power user!

\section{Revision Specifiers}\label{svn.tour.revs.specifiers}

As we described in \hyperref[svn.basic.in-action.revs]{Revisions},
revision numbers in Subversion are pretty straightforward---integers
that keep getting larger as you commit more changes to your versioned
data. Still, it doesn\textquotesingle t take long before you can no
longer remember exactly what happened in each and every revision.
Fortunately, the typical Subversion workflow doesn\textquotesingle t
often demand that you supply arbitrary revisions to the Subversion
operations you perform. For operations that \emph{do} require a revision
specifier, you generally supply a revision number that you saw in a
commit email, in the output of some other Subversion operation, or in
some other context that would give meaning to that particular number.

Referring to revision numbers with an ``\texttt{r}'' prefix
(\texttt{r314}, for example) is an established practice in Subversion
communities, and is both supported and encouraged by many
Subversion-related tools. In most places where you would specify a bare
revision number on the command line, you may also use the
\texttt{r}\textless NNN\textgreater{} syntax.

But occasionally, you need to pinpoint a moment in time for which you
don\textquotesingle t already have a revision number memorized or handy.
So besides the integer revision numbers, \texttt{svn} allows as input
some additional forms of revision specifiers: revision keywords and
revision dates.

The various forms of Subversion revision specifiers can be mixed and
matched when used to specify revision ranges. For example, you can use
\texttt{-r\ REV1:REV2} where \textless REV1\textgreater{} is a revision
keyword and \textless REV2\textgreater{} is a revision number, or where
\textless REV1\textgreater{} is a date and \textless REV2\textgreater{}
is a revision keyword, and so on. The individual revision specifiers are
independently evaluated, so you can put whatever you want on the
opposite sides of that colon.

\subsection{Revision Keywords}\label{svn.tour.revs.keywords}

The Subversion client understands a number of revision keywords. These
keywords can be used instead of integer arguments to the
\texttt{-\/-revision} (\texttt{-r}) option, and are resolved into
specific revision numbers by Subversion:

\begin{description}
\item[\texttt{HEAD}]
The latest (or ``youngest'') revision in the repository.
\item[\texttt{BASE}]
The revision number of an item in a working copy. If the item has been
locally modified, this refers to the way the item appears without those
local modifications.
\item[\texttt{COMMITTED}]
The most recent revision prior to, or equal to, \texttt{BASE}, in which
an item changed.
\item[\texttt{PREV}]
The revision immediately \emph{before} the last revision in which an
item changed. Technically, this boils down to \texttt{COMMITTED}-1.
\end{description}

As can be derived from their descriptions, the \texttt{PREV},
\texttt{BASE}, and \texttt{COMMITTED} revision keywords are used only
when referring to a working copy path---they don\textquotesingle t apply
to repository URLs. \texttt{HEAD}, on the other hand, can be used in
conjunction with both of these path types.

Here are some examples of revision keywords in action:

\begin{verbatim}
$ svn diff -r PREV:COMMITTED foo.c
# shows the last change committed to foo.c

$ svn log -r HEAD
# shows log message for the latest repository commit

$ svn diff -r HEAD
# compares your working copy (with all of its local changes) to the
# latest version of that tree in the repository

$ svn diff -r BASE:HEAD foo.c
# compares the unmodified version of foo.c with the latest version of
# foo.c in the repository

$ svn log -r BASE:HEAD
# shows all commit logs for the current versioned directory since you
# last updated

$ svn update -r PREV foo.c
# rewinds the last change on foo.c, decreasing foo.c's working revision

$ svn diff -r BASE:14 foo.c
# compares the unmodified version of foo.c with the way foo.c looked
# in revision 14
\end{verbatim}

\subsection{Revision Dates}\label{svn.tour.revs.dates}

Revision numbers reveal nothing about the world outside the version
control system, but sometimes you need to correlate a moment in real
time with a moment in version history. To facilitate this, the
\texttt{-\/-revision} (\texttt{-r}) option can also accept as input date
specifiers wrapped in curly braces (\texttt{\{} and \texttt{\}}).
Subversion accepts the standard ISO-8601 date and time formats, plus a
few others. Here are some examples.

\begin{verbatim}
$ svn update -r {2006-02-17}
$ svn update -r {15:30}
$ svn update -r {15:30:00.200000}
$ svn update -r {"2006-02-17 15:30"}
$ svn update -r {"2006-02-17 15:30 +0230"}
$ svn update -r {2006-02-17T15:30}
$ svn update -r {2006-02-17T15:30Z}
$ svn update -r {2006-02-17T15:30-04:00}
$ svn update -r {20060217T1530}
$ svn update -r {20060217T1530Z}
$ svn update -r {20060217T1530-0500}
…
\end{verbatim}

Keep in mind that most shells will require you to, at a minimum, quote
or otherwise escape any spaces that are included as part of revision
date specifiers. Certain shells may also take issue with the unescaped
use of curly braces, too. Consult your shell\textquotesingle s
documentation for the requirements specific to your environment.

When you specify a date, Subversion resolves that date to the most
recent revision of the repository as of that date, and then continues to
operate against that resolved revision number:

\begin{verbatim}
$ svn log -r {2006-11-28}
------------------------------------------------------------------------
r12 | ira | 2006-11-27 12:31:51 -0600 (Mon, 27 Nov 2006) | 6 lines
…
\end{verbatim}

Is Subversion a Day Early?

If you specify a single date as a revision without specifying a time of
day (for example \texttt{2006-11-27}), you may think that Subversion
should give you the last revision that took place on the 27th of
November. Instead, you\textquotesingle ll get back a revision from the
26th, or even earlier. Remember that Subversion will find the \emph{most
recent revision of the repository} as of the date you give. If you give
a date without a timestamp, such as \texttt{2006-11-27}, Subversion
assumes a time of 00:00:00, so looking for the most recent revision
won\textquotesingle t return anything on the 27th.

If you want to include the 27th in your search, you can either specify
the 27th with the time (\texttt{\{"2006-11-27\ 23:59"\}}), or just
specify the next day (\texttt{\{2006-11-28\}}).

You can also use a range of dates. Subversion will find all revisions
between both dates, inclusive:

\begin{verbatim}
$ svn log -r {2006-11-20}:{2006-11-29}
…
\end{verbatim}

Subversion\textquotesingle s ability to correctly convert revision dates
into real revision numbers depends on revision datestamps maintaining a
sequential ordering---the younger the revision, the younger its
datestamp. But datestamps are stored in the unversioned, modifiable
\texttt{svn:date} property of the revision (see
\hyperref[svn.advanced.props]{Properties}), so it is possible for
revision datestamps to get out of sequence. Now, most of
Subversion\textquotesingle s operations are unaffected by this
situation---after all, the revision number itself is the primary
identifier of each revision. But if the datestamp ordering
isn\textquotesingle t maintained, you will likely find that trying to
use dates to specify revision ranges in your repository
doesn\textquotesingle t always return the data you might have expected.
Combining the histories of multiple repositories into a single one (as
described in \hyperref[svn.reposadmin.maint.migrate]{Migrating
Repository Data Elsewhere}) is the most common cause of this scenario.

\section{Peg and Operative Revisions}\label{svn.advanced.pegrevs}

We copy, move, rename, and completely replace files and directories on
our computers all the time. And your version control system
shouldn\textquotesingle t get in the way of your doing these things with
your version-controlled files and directories, either.
Subversion\textquotesingle s file management support is quite
liberating, affording almost as much flexibility for versioned files as
you\textquotesingle d expect when manipulating your unversioned ones.
But that flexibility means that across the lifetime of your repository,
a given versioned object might have many paths, and a given path might
represent several entirely different versioned objects. This introduces
a certain level of complexity to your interactions with those paths and
objects.

Subversion is pretty smart about noticing when an
object\textquotesingle s version history includes such ``changes of
address.'' For example, if you ask for the revision history log of a
particular file that was renamed last week, Subversion happily provides
all those logs---the revision in which the rename itself happened, plus
the logs of relevant revisions both before and after that rename. So,
most of the time, you don\textquotesingle t even have to think about
such things. But occasionally, Subversion needs your help to clear up
ambiguities.

The simplest example of this occurs when a directory or file is deleted
from version control, and then a new directory or file is created with
the same name and added to version control. The thing you deleted and
the thing you later added aren\textquotesingle t the same thing. They
merely happen to have had the same path---\texttt{/trunk/object}, for
example. What, then, does it mean to ask Subversion about the history of
\texttt{/trunk/object}? Are you asking about the thing currently at that
location, or the old thing you deleted from that location? Are you
asking about the operations that have happened to \emph{all} the objects
that have ever lived at that path? Subversion needs a hint about what
you really want.

And thanks to moves, versioned object history can get far more twisted
than even that. For example, you might have a directory named
\texttt{concept}, containing some nascent software project
you\textquotesingle ve been toying with. Eventually, though, that
project matures to the point that the idea seems to actually have some
wings, so you do the unthinkable and decide to give the project a
name.\footnote{``You\textquotesingle re not supposed to name it. Once
  you name it, you start getting attached to it.''---Mike Wazowski}
Let\textquotesingle s say you called your software Frabnaggilywort. At
this point, it makes sense to rename the directory to reflect the
project\textquotesingle s new name, so \texttt{concept} is renamed to
\texttt{frabnaggilywort}. Life goes on, Frabnaggilywort releases a 1.0
version and is downloaded and used daily by hordes of people aiming to
improve their lives.

It\textquotesingle s a nice story, really, but it
doesn\textquotesingle t end there. Entrepreneur that you are,
you\textquotesingle ve already got another think in the tank. So you
make a new directory, \texttt{concept}, and the cycle begins again. In
fact, the cycle begins again many times over the years, each time
starting with that old \texttt{concept} directory, then sometimes seeing
that directory renamed as the idea cures, sometimes seeing it deleted
when you scrap the idea. Or, to get really sick, maybe you rename
\texttt{concept} to something else for a while, but later rename the
thing back to \texttt{concept} for some reason.

In scenarios like these, attempting to instruct Subversion to work with
these reused paths can be a little like instructing a motorist in
Chicago\textquotesingle s West Suburbs to drive east down Roosevelt Road
and turn left onto Main Street. In a mere 20 minutes, you can cross
``Main Street'' in Wheaton, Glen Ellyn, and Lombard. And no, they
aren\textquotesingle t the same street. Our motorist---and our
Subversion---need a little more detail to do the right thing.

Fortunately, Subversion allows you to tell it exactly which Main Street
you meant. The mechanism used is called a peg revision, and you provide
these to Subversion for the sole purpose of identifying unique lines of
history. Because at most one versioned object may occupy a path at any
given time---or, more precisely, in any one revision---the combination
of a path and a peg revision is all that is needed to unambiguously
identify a specific line of history. Peg revisions are specified to the
Subversion command-line client using at syntax, so called because the
syntax involves appending an ``at sign'' (\texttt{@}) and the peg
revision to the end of the path with which the revision is associated.

But what of the \texttt{-\/-revision} (\texttt{-r}) of which
we\textquotesingle ve spoken so much in this book? That revision (or set
of revisions) is called the operative revision (or operative revision
range). Once a particular line of history has been identified using a
path and peg revision, Subversion performs the requested operation using
the operative revision(s). To map this to our Chicagoland streets
analogy, if we are told to go to 606 N. Main Street in
Wheaton,\footnote{606 N. Main Street, Wheaton, Illinois, is the home of
  the Wheaton \emph{History} Center. It seemed appropriate\ldots{}} we
can think of ``Main Street'' as our path and ``Wheaton'' as our peg
revision. These two pieces of information identify a unique path that
can be traveled (north or south on Main Street), and they keep us from
traveling up and down the wrong Main Street in search of our
destination. Now we throw in ``606 N.'' as our operative revision of
sorts, and we know \emph{exactly} where to go.

\protect\phantomsection\label{svn.advanced.pegrevs.algorithm}
The Peg Revision Algorithm

The Subversion command-line client performs the peg revision algorithm
any time it needs to resolve possible ambiguities in the paths and
revisions provided to it. Here\textquotesingle s an example of such an
invocation:

\begin{verbatim}
$ svn command -r OPERATIVE-REV item@PEG-REV
\end{verbatim}

If \textless OPERATIVE-REV\textgreater{} is older than
\textless PEG-REV\textgreater, the algorithm is as follows:

\begin{enumerate}
\def\labelenumi{\arabic{enumi}.}
\item
  Locate \textless item\textgreater{} in the revision identified by
  \textless PEG-REV\textgreater. There can be only one such object.
\item
  Trace the object\textquotesingle s history backwards (through any
  possible renames) to its ancestor in the revision
  \textless OPERATIVE-REV\textgreater.
\item
  Perform the requested action on that ancestor, wherever it is located,
  or whatever its name might be or might have been at that time.
\end{enumerate}

But what if \textless OPERATIVE-REV\textgreater{} is \emph{younger} than
\textless PEG-REV\textgreater? Well, that adds some complexity to the
theoretical problem of locating the path in
\textless OPERATIVE-REV\textgreater, because the path\textquotesingle s
history could have forked multiple times (thanks to copy operations)
between \textless PEG-REV\textgreater{} and
\textless OPERATIVE-REV\textgreater. And that\textquotesingle s not
all---Subversion doesn\textquotesingle t store enough information to
performantly trace an object\textquotesingle s history forward, anyway.
So the algorithm is a little different:

\begin{enumerate}
\def\labelenumi{\arabic{enumi}.}
\item
  Locate \textless item\textgreater{} in the revision identified by
  \textless OPERATIVE-REV\textgreater. There can be only one such
  object.
\item
  Trace the object\textquotesingle s history backward (through any
  possible renames) to its ancestor in the revision
  \textless PEG-REV\textgreater.
\item
  Verify that the object\textquotesingle s location (path-wise) in
  \textless PEG-REV\textgreater{} is the same as it is in
  \textless OPERATIVE-REV\textgreater. If that\textquotesingle s the
  case, at least the two locations are known to be directly related, so
  perform the requested action on the location in
  \textless OPERATIVE-REV\textgreater. Otherwise, relatedness was not
  established, so error out with a loud complaint that no viable
  location was found. (Someday, we expect that Subversion will be able
  to handle this usage scenario with more flexibility and grace.)
\end{enumerate}

Note that even when you don\textquotesingle t explicitly supply a peg
revision or operative revision, they are still present. For your
convenience, the default peg revision is \texttt{BASE} for working copy
items and \texttt{HEAD} for repository URLs. And when no operative
revision is provided, it defaults to being the same revision as the peg
revision.

Say that long ago we created our repository, and in revision 1 we added
our first \texttt{concept} directory, plus an \texttt{IDEA} file in that
directory talking about the concept. After several revisions in which
real code was added and tweaked, we, in revision 20, renamed this
directory to \texttt{frabnaggilywort}. By revision 27, we had a new
concept, a new \texttt{concept} directory to hold it, and a new
\texttt{IDEA} file to describe it. And then five years and thousands of
revisions flew by, just like they would in any good romance story.

Now, years later, we wonder what the \texttt{IDEA} file looked like back
in revision 1. But Subversion needs to know whether we are asking about
how the \emph{current} file looked back in revision 1, or whether we are
asking for the contents of whatever file lived at \texttt{concept/IDEA}
in revision 1. Certainly those questions have different answers, and
because of peg revisions, you can ask those questions. To find out how
the current \texttt{IDEA} file looked in that old revision, you run:

\begin{verbatim}
$ svn cat -r 1 concept/IDEA 
svn: E195012: Unable to find repository location for 'concept/IDEA' in revision 1
\end{verbatim}

Of course, in this example, the current \texttt{IDEA} file
didn\textquotesingle t exist yet in revision 1, so Subversion gives an
error. The previous command is shorthand for a longer notation which
explicitly lists a peg revision. The expanded notation is:

\begin{verbatim}
$ svn cat -r 1 concept/IDEA@BASE
svn: E195012: Unable to find repository location for 'concept/IDEA' in revision 1
\end{verbatim}

And when executed, it has the expected results.

The perceptive reader is probably wondering at this point whether the
peg revision syntax causes problems for working copy paths or URLs that
actually have at signs in them. After all, how does \texttt{svn} know
whether \texttt{news@11} is the name of a directory in my tree or just a
syntax for ``revision 11 of \texttt{news}''? Thankfully, while
\texttt{svn} will always assume the latter, there is a trivial
workaround. You need only append an at sign to the end of the path, such
as \texttt{news@11@}. \texttt{svn} cares only about the last at sign in
the argument, and it is not considered illegal to omit a literal peg
revision specifier after that at sign. This workaround even applies to
paths that end in an at sign---you would use \texttt{filename@@} to talk
about a file named \texttt{filename@}.

Let\textquotesingle s ask the other question, then---in revision 1, what
were the contents of whatever file occupied the address
\texttt{concepts/IDEA} at the time? We\textquotesingle ll use an
explicit peg revision to help us out.

\begin{verbatim}
$ svn cat concept/IDEA@1
The idea behind this project is to come up with a piece of software
that can frab a naggily wort.  Frabbing naggily worts is tricky
business, and doing it incorrectly can have serious ramifications, so
we need to employ over-the-top input validation and data verification
mechanisms.
\end{verbatim}

Notice that we didn\textquotesingle t provide an operative revision this
time. That\textquotesingle s because when no operative revision is
specified, Subversion assumes a default operative revision
that\textquotesingle s the same as the peg revision.

As you can see, the output from our operation appears to be correct. The
text even mentions frabbing naggily worts, so this is almost certainly
the file that describes the software now called Frabnaggilywort. In
fact, we can verify this using the combination of an explicit peg
revision and explicit operative revision. We know that in \texttt{HEAD},
the Frabnaggilywort project is located in the \texttt{frabnaggilywort}
directory. So we specify that we want to see how the line of history
identified in \texttt{HEAD} as the path \texttt{frabnaggilywort/IDEA}
looked in revision 1.

\begin{verbatim}
$ svn cat -r 1 frabnaggilywort/IDEA@HEAD
The idea behind this project is to come up with a piece of software
that can frab a naggily wort.  Frabbing naggily worts is tricky
business, and doing it incorrectly can have serious ramifications, so
we need to employ over-the-top input validation and data verification
mechanisms.
\end{verbatim}

And the peg and operative revisions need not be so trivial, either. For
example, say \texttt{frabnaggilywort} had been deleted from
\texttt{HEAD}, but we know it existed in revision 20, and we want to see
the diffs for its \texttt{IDEA} file between revisions 4 and 10. We can
use peg revision 20 in conjunction with the URL that would have held
Frabnaggilywort\textquotesingle s \texttt{IDEA} file in revision 20, and
then use 4 and 10 as our operative revision range.

\begin{verbatim}
$ svn diff -r 4:10 http://svn.red-bean.com/projects/frabnaggilywort/IDEA@20
Index: frabnaggilywort/IDEA
===================================================================
--- frabnaggilywort/IDEA    (revision 4)
+++ frabnaggilywort/IDEA    (revision 10)
@@ -1,5 +1,5 @@
-The idea behind this project is to come up with a piece of software
-that can frab a naggily wort.  Frabbing naggily worts is tricky
-business, and doing it incorrectly can have serious ramifications, so
-we need to employ over-the-top input validation and data verification
-mechanisms.
+The idea behind this project is to come up with a piece of
+client-server software that can remotely frab a naggily wort.
+Frabbing naggily worts is tricky business, and doing it incorrectly
+can have serious ramifications, so we need to employ over-the-top
+input validation and data verification mechanisms.
\end{verbatim}

Fortunately, most folks aren\textquotesingle t faced with such complex
situations. But when you are, remember that peg revisions are that extra
hint Subversion needs to clear up ambiguity.

\section{Properties}\label{svn.advanced.props}

We\textquotesingle ve already covered in detail how Subversion stores
and retrieves various versions of files and directories in its
repository. Whole chapters have been devoted to this most fundamental
piece of functionality provided by the tool. And if the versioning
support stopped there, Subversion would still be complete from a version
control perspective.

But it doesn\textquotesingle t stop there.

In addition to versioning your directories and files, Subversion
provides interfaces for adding, modifying, and removing versioned
metadata on each of your versioned directories and files. We refer to
this metadata as properties, and they can be thought of as two-column
tables that map property names to arbitrary values attached to each item
in your working copy. Generally speaking, the names and values of the
properties can be whatever you want them to be, with the constraint that
the names must contain only ASCII characters. And the best part about
these properties is that they, too, are versioned, just like the textual
contents of your files. You can modify, commit, and revert property
changes as easily as you can file content changes. And the sending and
receiving of property changes occurs as part of your typical commit and
update operations---you don\textquotesingle t have to change your basic
processes to accommodate them.

Subversion has reserved the set of properties whose names begin with
\texttt{svn:} as its own. While there are only a handful of such
properties in use today, you should avoid creating custom properties for
your own needs whose names begin with this prefix. Otherwise, you run
the risk that a future release of Subversion will grow support for a
feature or behavior driven by a property of the same name but with
perhaps an entirely different interpretation.

Properties show up elsewhere in Subversion, too. Just as files and
directories may have arbitrary property names and values attached to
them, each revision as a whole may have arbitrary properties attached to
it. The same constraints apply---human-readable names and
anything-you-want binary values. The main difference is that revision
properties are not versioned. In other words, if you change the value
of, or delete, a revision property, there\textquotesingle s no way,
within the scope of Subversion\textquotesingle s functionality, to
recover the previous value.

Subversion has no particular policy regarding the use of properties. It
asks only that you do not use property names that begin with the prefix
\texttt{svn:} as that\textquotesingle s the namespace that it sets aside
for its own use. And Subversion does, in fact, use properties---both the
versioned and unversioned variety. Certain versioned properties have
special meaning or effects when found on files and directories, or they
house a particular bit of information about the revisions on which they
are found. Certain revision properties are automatically attached to
revisions by Subversion\textquotesingle s commit process, and they carry
information about the revision. Most of these properties are mentioned
elsewhere in this or other chapters as part of the more general topics
to which they are related. For an exhaustive list of
Subversion\textquotesingle s predefined properties, see
\hyperref[svn.advanced.props.ref]{Subversion\textquotesingle s Reserved
Properties}.

While Subversion automatically attaches properties (\texttt{svn:date},
\texttt{svn:author}, \texttt{svn:log}, and so on) to revisions, it does
\emph{not} presume thereafter the existence of those properties, and
neither should you or the tools you use to interact with your
repository. Revision properties can be deleted programmatically or via
the client (if allowed by the repository hooks) without damaging
Subversion\textquotesingle s ability to function. So, when writing
scripts which operate on your Subversion repository data, do not make
the mistake of assuming that any particular revision property exists on
a revision.

In this section, we will examine the utility---both to users of
Subversion and to Subversion itself---of property support.
You\textquotesingle ll learn about the property-related \texttt{svn}
subcommands and how property modifications affect your normal Subversion
workflow.

\subsection{Why Properties?}\label{svn.advanced.props.why}

Just as Subversion uses properties to store extra information about the
files, directories, and revisions that it contains, you might also find
properties to be of similar use. You might find it useful to have a
place close to your versioned data to hang custom metadata about that
data.

Say you wish to design a web site that houses many digital photos and
displays them with captions and a datestamp. Now, your set of photos is
constantly changing, so you\textquotesingle d like to have as much of
this site automated as possible. These photos can be quite large, so as
is common with sites of this nature, you want to provide smaller
thumbnail images to your site visitors.

Now, you can get this functionality using traditional files. That is,
you can have your \texttt{image123.jpg} and an
\texttt{image123-thumbnail.jpg} side by side in a directory. Or if you
want to keep the filenames the same, you might have your thumbnails in a
different directory, such as \texttt{thumbnails/image123.jpg}. You can
also store your captions and datestamps in a similar fashion, again
separated from the original image file. But the problem here is that
your collection of files multiplies with each new photo added to the
site.

Now consider the same web site deployed in a way that makes use of
Subversion\textquotesingle s file properties. Imagine having a single
image file, \texttt{image123.jpg}, with properties set on that file that
are named \texttt{caption}, \texttt{datestamp}, and even
\texttt{thumbnail}. Now your working copy directory looks much more
manageable---in fact, it looks to the casual browser like there are
nothing but image files in it. But your automation scripts know better.
They know that they can use \texttt{svn} (or better yet, they can use
the Subversion language bindings---see
\hyperref[svn.developer.usingapi]{Using the APIs}) to dig out the extra
information that your site needs to display without having to read an
index file or play path manipulation games.

While Subversion places few restrictions on the names and values you use
for properties, it has not been designed to optimally carry large
property values or large sets of properties on a given file or
directory. Subversion commonly holds all the property names and values
associated with a single item in memory at the same time, which can
cause detrimental performance or failed operations when extremely large
property sets are used.

Custom revision properties are also frequently used. One common such use
is a property whose value contains an issue tracker ID with which the
revision is associated, perhaps because the change made in that revision
fixes a bug filed in the tracker issue with that ID. Other uses include
hanging more friendly names on the revision---it might be hard to
remember that revision 1935 was a fully tested revision. But if
there\textquotesingle s, say, a \texttt{test-results} property on that
revision with the value \texttt{all\ passing}, that\textquotesingle s
meaningful information to have. And Subversion allows you to easily do
this via the \texttt{-\/-with-revprop} option of the
\texttt{svn\ commit} command:

\begin{verbatim}
$ svn commit -m "Fix up the last remaining known regression bug." \
             --with-revprop "test-results=all passing"
Sending        lib/crit_bits.c
Transmitting file data .
Committed revision 912.
$
\end{verbatim}

Searchability (or, Why \emph{Not} Properties)

For all their utility, Subversion properties---or, more accurately, the
available interfaces to them---have a major shortcoming: while it is a
simple matter to \emph{set} a custom property, \emph{finding} that
property later is a whole different ball of wax.

Trying to locate a custom revision property generally involves
performing a linear walk across all the revisions of the repository,
asking of each revision, ``Do you have the property I\textquotesingle m
looking for?'' Use the \texttt{-\/-with-all-revprops} option with the
\texttt{svn\ log} command\textquotesingle s XML output mode to
facilitate this search. Notice the presence of the custom revision
property \texttt{testresults} in the following output:

\begin{verbatim}
$ svn log --with-all-revprops --xml lib/crit_bits.c
<?xml version="1.0"?>
<log>
<logentry
   revision="912">
<author>harry</author>
<date>2011-07-29T14:47:41.169894Z</date>
<msg>Fix up the last remaining known regression bug.</msg>
<revprops>
<property
   name="testresults">all passing</property>
</revprops>
</logentry>
…
$
\end{verbatim}

Trying to find a custom versioned property is painful, too, and often
involves a recursive \texttt{svn\ propget} across an entire working
copy. In your situation, that might not be as bad as a linear walk
across all revisions. But it certainly leaves much to be desired in
terms of both performance and likelihood of success, especially if the
scope of your search would require a working copy from the root of your
repository.

For this reason, you might choose---especially in the revision property
use case---to simply add your metadata to the revision\textquotesingle s
log message using some policy-driven (and perhaps programmatically
enforced) formatting that is designed to be quickly parsed from the
output of \texttt{svn\ log}. It is quite common to see the following in
Subversion log messages:

\begin{verbatim}
Issue(s): IZ2376, IZ1919
Reviewed by:  sally

This fixes a nasty segfault in the wort frabbing process
…
\end{verbatim}

But here again lies some misfortune. Subversion doesn\textquotesingle t
yet provide a log message templating mechanism, which would go a long
way toward helping users be consistent with the formatting of their
log-embedded revision metadata.

\subsection{Manipulating Properties}\label{svn.advanced.props.manip}

The \texttt{svn} program affords a few ways to add or modify file and
directory properties. For properties with short, human-readable values,
perhaps the simplest way to add a new property is to specify the
property name and value on the command line of the \texttt{svn\ propset}
subcommand:

\begin{verbatim}
$ svn propset copyright '(c) 2006 Red-Bean Software' calc/button.c
property 'copyright' set on 'calc/button.c'
$
\end{verbatim}

But we\textquotesingle ve been touting the flexibility that Subversion
offers for your property values. And if you are planning to have a
multiline textual, or even binary, property value, you probably do not
want to supply that value on the command line. So the
\texttt{svn\ propset} subcommand takes a \texttt{-\/-file} (\texttt{-F})
option for specifying the name of a file that contains the new property
value.

\begin{verbatim}
$ svn propset license -F /path/to/LICENSE calc/button.c
property 'license' set on 'calc/button.c'
$
\end{verbatim}

There are some restrictions on the names you can use for properties. A
property name must start with a letter, a colon (\texttt{:}), or an
underscore (\texttt{\_}); after that, you can also use digits, hyphens
(\texttt{-}), and periods (\texttt{.}).\footnote{If
  you\textquotesingle re familiar with XML, this is pretty much the
  ASCII subset of the syntax for XML ``Name''.}

In addition to the \texttt{propset} command, the \texttt{svn} program
supplies the \texttt{propedit} command. This command uses the configured
editor program (see \hyperref[svn.advanced.confarea.opts.config]{General
configuration}) to add or modify properties. When you run the command,
\texttt{svn} invokes your editor program on a temporary file that
contains the current value of the property (or that is empty, if you are
adding a new property). Then, you just modify that value in your editor
program until it represents the new value you wish to store for the
property, save the temporary file, and then exit the editor program. If
Subversion detects that you\textquotesingle ve actually changed the
existing value of the property, it will accept that as the new property
value. If you exit your editor without making any changes, no property
modification will occur:

\begin{verbatim}
$ svn propedit copyright calc/button.c  ### exit the editor without changes
No changes to property 'copyright' on 'calc/button.c'
$
\end{verbatim}

We should note that, as with other \texttt{svn} subcommands, those
related to properties can act on multiple paths at once. This enables
you to modify properties on whole sets of files with a single command.
For example, we could have done the following:

\begin{verbatim}
$ svn propset copyright '(c) 2006 Red-Bean Software' calc/*
property 'copyright' set on 'calc/Makefile'
property 'copyright' set on 'calc/button.c'
property 'copyright' set on 'calc/integer.c'
…
$
\end{verbatim}

All of this property adding and editing isn\textquotesingle t really
very useful if you can\textquotesingle t easily get the stored property
value. So the \texttt{svn} program supplies two subcommands for
displaying the names and values of properties stored on files and
directories. The \texttt{svn\ proplist} command will list the names of
properties that exist on a path. Once you know the names of the
properties on the node, you can request their values individually using
\texttt{svn\ propget}. This command will, given a property name and a
path (or set of paths), print the value of the property to the standard
output stream.

\begin{verbatim}
$ svn proplist calc/button.c
Properties on 'calc/button.c':
  copyright
  license
$ svn propget copyright calc/button.c
(c) 2006 Red-Bean Software
\end{verbatim}

There\textquotesingle s even a variation of the \texttt{proplist}
command that will list both the name and the value for all of the
properties. Simply supply the \texttt{-\/-verbose} (\texttt{-v}) option.

\begin{verbatim}
$ svn proplist -v calc/button.c
Properties on 'calc/button.c':
  copyright
    (c) 2006 Red-Bean Software
  license
    ================================================================
    Copyright (c) 2006 Red-Bean Software.  All rights reserved.

    Redistribution and use in source and binary forms, with or without
    modification, are permitted provided that the following conditions 
    are met:

    1. Redistributions of source code must retain the above copyright
    notice, this list of conditions, and the recipe for Fitz's famous
    red-beans-and-rice.
    …
\end{verbatim}

The last property-related subcommand is \texttt{propdel}. Since
Subversion allows you to store properties with empty values, you
can\textquotesingle t remove a property altogether using
\texttt{svn\ propedit} or \texttt{svn\ propset}. For example, this
command will \emph{not} yield the desired effect:

\begin{verbatim}
$ svn propset license "" calc/button.c
property 'license' set on 'calc/button.c'
$ svn proplist -v calc/button.c
Properties on 'calc/button.c':
  copyright
    (c) 2006 Red-Bean Software
  license
    
$
\end{verbatim}

You need to use the \texttt{propdel} subcommand to delete properties
altogether. The syntax is similar to the other property commands:

\begin{verbatim}
$ svn propdel license calc/button.c
property 'license' deleted from 'calc/button.c'.
$ svn proplist -v calc/button.c
Properties on 'calc/button.c':
  copyright
    (c) 2006 Red-Bean Software
$
\end{verbatim}

Remember those unversioned revision properties? You can modify those,
too, using the same \texttt{svn} subcommands that we just described.
Simply add the \texttt{-\/-revprop} command-line parameter and specify
the revision whose property you wish to modify. Since revisions are
global, you don\textquotesingle t need to specify a target path to these
property-related commands so long as you are positioned in a working
copy of the repository whose revision property you wish to modify.
Otherwise, you can simply provide the URL of any path in the repository
of interest (including the repository\textquotesingle s root URL). For
example, you might want to replace the commit log message of an existing
revision.\footnote{Fixing spelling errors, grammatical gotchas, and
  ``just-plain-wrongness'' in commit log messages is perhaps the most
  common use case for the \texttt{-\/-revprop} option.} If your current
working directory is part of a working copy of your repository, you can
simply run the \texttt{svn\ propset} command with no target path:

\begin{verbatim}
$ svn propset svn:log "* button.c: Fix a compiler warning." -r11 --revprop
property 'svn:log' set on repository revision '11'
$
\end{verbatim}

But even if you haven\textquotesingle t checked out a working copy from
that repository, you can still effect the property change by providing
the repository\textquotesingle s root URL:

\begin{verbatim}
$ svn propset svn:log "* button.c: Fix a compiler warning." -r11 --revprop \
              http://svn.example.com/repos/project
property 'svn:log' set on repository revision '11'
$
\end{verbatim}

Note that the ability to modify these unversioned properties must be
explicitly added by the repository administrator (see
\hyperref[svn.reposadmin.maint.setlog]{Commit Log Message Correction}).
That\textquotesingle s because the properties aren\textquotesingle t
versioned, so you run the risk of losing information if you
aren\textquotesingle t careful with your edits. The repository
administrator can set up methods to protect against this loss, and by
default, modification of unversioned properties is disabled.

Users should, where possible, use \texttt{svn\ propedit} instead of
\texttt{svn\ propset}. While the end result of the commands is
identical, the former will allow them to see the current value of the
property that they are about to change, which helps them to verify that
they are, in fact, making the change they think they are making. This is
especially true when modifying unversioned revision properties. Also, it
is significantly easier to modify multiline property values in a text
editor than at the command line.

\subsection{Properties and the Subversion
Workflow}\label{svn.advanced.props.workflow}

Now that you are familiar with all of the property-related \texttt{svn}
subcommands, let\textquotesingle s see how property modifications affect
the usual Subversion workflow. As we mentioned earlier, file and
directory properties are versioned, just like your file contents. As a
result, Subversion provides the same opportunities for merging---cleanly
or with conflicts---someone else\textquotesingle s modifications into
your own.

As with file contents, your property changes are local modifications,
made permanent only when you commit them to the repository with
\texttt{svn\ commit}. Your property changes can be easily unmade,
too---the \texttt{svn\ revert} command will restore your files and
directories to their unedited states---contents, properties, and all.
Also, you can receive interesting information about the state of your
file and directory properties by using the \texttt{svn\ status} and
\texttt{svn\ diff} commands.

\begin{verbatim}
$ svn status calc/button.c
 M      calc/button.c
$ svn diff calc/button.c
Property changes on: calc/button.c
___________________________________________________________________
Added: copyright
## -0,0 +1 ##
+(c) 2006 Red-Bean Software
$
\end{verbatim}

Notice how the \texttt{status} subcommand displays \texttt{M} in the
second column instead of the first. That is because we have modified the
properties on \texttt{calc/button.c}, but not its textual contents. Had
we changed both, we would have seen \texttt{M} in the first column, too.
(We cover \texttt{svn\ status} in
\hyperref[svn.tour.cycle.examine.status]{See an overview of your
changes}).

Property Conflicts

As with file contents, local property modifications can conflict with
changes committed by someone else. If you update your working copy
directory and receive property changes on a versioned object that clash
with your own, Subversion will report that the object is in a conflicted
state.

\begin{verbatim}
$ svn update calc
Updating 'calc':
M  calc/Makefile.in
Conflict for property 'linecount' discovered on 'calc/button.c'.
Select: (p) postpone, (df) diff-full, (e) edit,
        (s) show all options: p
 C calc/button.c
Updated to revision 143.
Summary of conflicts:
  Property conflicts: 1
$ 
\end{verbatim}

Subversion will also create, in the same directory as the conflicted
object, a file with a \texttt{.prej} extension that contains the details
of the conflict. You should examine the contents of this file so you can
decide how to resolve the conflict. Until the conflict is resolved, you
will see a \texttt{C} in the second column of \texttt{svn\ status}
output for that object, and attempts to commit your local modifications
will fail.

\begin{verbatim}
$ svn status calc
 C      calc/button.c
?       calc/button.c.prej
$ cat calc/button.c.prej 
Trying to change property 'linecount' from '1267' to '1301',
but property has been locally changed from '1267' to '1256'.
$
\end{verbatim}

To resolve property conflicts, simply ensure that the conflicting
properties contain the values that they should, and then use the
\texttt{svn\ resolve\ -\/-accept=working} command to alert Subversion
that you have manually resolved the problem.

You might also have noticed the nonstandard way that Subversion
currently displays property differences. You can still use
\texttt{svn\ diff} and redirect its output to create a usable patch
file. The \texttt{patch} program will ignore property patches---as a
rule, it ignores any noise it can\textquotesingle t understand. This
does, unfortunately, mean that to fully apply a patch generated by
\texttt{svn\ diff} using \texttt{patch}, any property modifications will
need to be applied by hand.

Subversion 1.7 improves this situation in two ways. First, its
nonstandard display of property differences is at least
machine-readable---an improvement over the display of properties in
versions prior to 1.7. But Subversion 1.7 also introduces the
\texttt{svn\ patch} subcommand, designed specifically to handle the
additional information which \texttt{svn\ diff}\textquotesingle s output
can carry, applying those changes to the Subversion working copy. Of
specific relevance to our topic, property differences present in patch
files generated by \texttt{svn\ diff} in Subversion 1.7 or better can be
automatically applied to a working copy by the \texttt{svn\ patch}
command. For more about \texttt{svn\ patch}, see
\hyperref[svn.ref.svn.c.patch]{svn patch} in \hyperref[svn.ref.svn]{svn
Reference---Subversion Command-Line Client}.

There\textquotesingle s one exception to how property changes are
reported by \texttt{svn\ diff}: changes to Subversion\textquotesingle s
special \texttt{svn:mergeinfo} property---used to track information
about merges which have been performed in your repository---are
described in a more human-readable fashion. This is quite helpful to the
humans who have to read those descriptions. But it also serves to cause
patching programs (including \texttt{svn\ patch}) to skip those change
descriptions as noise. This might sound like a bug, but it really
isn\textquotesingle t because this property is intended to be managed
solely by the \texttt{svn\ merge} subcommand. For more about merge
tracking, see \hyperref[svn.branchmerge]{Branching and Merging}.

\subsection{Inherited Properties}\label{svn.advanced.props.inheritable}

Subversion 1.8 introduces the concept of inherited properties. There is
really nothing special about a property that makes it inheritable. In
fact, all versioned properties are inheritable! The main difference
between versioned properties before 1.8 and after is that the latter
provides a mechanism to find the properties set on a target
path\textquotesingle s \emph{parents}, even if those parents are not
found within the working copy.

Generic property inheritance manifests itself in a few commands. First,
the \texttt{svn\ proplist} and \texttt{svn\ propget} subcommands can
retrieve all the properties on a URL\textquotesingle s or a working copy
path\textquotesingle s parents by using the
\texttt{-\/-show-inherited-props} option. You might think of this as the
opposite of a \texttt{-\/-recursive} subcommand operation---instead of
recursing "down" into a target\textquotesingle s subdirectories,
subcommands with the \texttt{-\/-show-inherited-props} option look "up"
into the target\textquotesingle s parent directories. The
\texttt{svnlook\ propget} and \texttt{svnlook\ proplist} subcommands
also use the \texttt{-\/-show-inherited-props} option in a similar
fashion.

Let\textquotesingle s look at an example of how this works. The
following recursive propget on the root of our working copy finds that
the \texttt{svn:auto-props} property is set on both the target of the
subcommand and one of its subdirectories \texttt{site}:

\begin{verbatim}
$ svn pg svn:auto-props --verbose -R .
Properties on '.':
  svn:auto-props
    *.py = svn:eol-style=native
    *.c = svn:eol-style=native
    *.h = svn:eol-style=native

Properties on 'site':
  svn:auto-props
    *.html = svn:eol-style=native
\end{verbatim}

If we were to instead make the target of the subcommand the subdirectory
\texttt{site}, then using the \texttt{-\/-show-inherited-props} option,
we find that the \texttt{svn:auto-props} property is set on the target
\emph{and} its parent. The parent\textquotesingle s properties are
called out as "inherited":

\begin{verbatim}
$ svn pg svn:auto-props --verbose --show-inherited-props site
Inherited properties on 'site',
from '.':
  svn:auto-props
    *.py = svn:eol-style=native
    *.c = svn:eol-style=native
    *.h = svn:eol-style=native

Properties on 'site':
  svn:auto-props
    *.html = svn:eol-style=native
\end{verbatim}

In the prior examples the root of the working copy corresponds to the
root of the repository, but properties can also be inherited from
outside the working copy when this is not the case.
Let\textquotesingle s checkout the \texttt{site} directory from the
prior example, making it the root of our working copy:

\begin{verbatim}
$ svn co http://svn.example.com/repos site-wc
A    site-wc/publish
A    site-wc/publish/ch2.html
A    site-wc/publish/news.html
A    site-wc/publish/ch3.html
A    site-wc/publish/faq.html
A    site-wc/publish/index.html
A    site-wc/publish/ch1.html
 U   site-wc
Checked out revision 19.

$ cd site-wc
\end{verbatim}

Now when we check for inherited properties on a working copy path we can
see that one property is inherited from a working copy parent and one
from a repository parent representing a location "above" the root of the
working copy:

\begin{verbatim}
$ svn pg svn:auto-props --verbose --show-inherited-props publish
Inherited properties on 'publish',
from 'http://svn.example.com/repos':
  svn:auto-props
    *.py = svn:eol-style=native
    *.c = svn:eol-style=native
    *.h = svn:eol-style=native

Inherited properties on 'publish',
from '.':
  svn:auto-props
    *.html = svn:eol-style=native
\end{verbatim}

You can only inherit properties from repository paths which you have
read authorization to---see
\hyperref[svn.serverconfig.svnserve.auth]{Built-in Authentication and
Authorization} and \hyperref[svn.serverconfig.httpd.authz]{Authorization
Options}. If you don\textquotesingle t have read authorization to a
parent path then it will appear as if the parent has no properties set
on it.

As mentioned above, the \texttt{svnlook\ proplist} and
\texttt{svnlook\ propget} commands also support the
\texttt{-\/-show-inherited-props} option, but instead of reporting the
inherited props by working copy path or URL, they are listed by
repository paths:

\begin{verbatim}
$ svnlook pg repos svn:auto-props /site/publish --show-inherited-props -v
Inherited properties on '/site/publish',
from '/':
  svn:auto-props
    *.py = svn:eol-style=native
    *.c = svn:eol-style=native
    *.h = svn:eol-style=native

Inherited properties on '/site/publish',
from '/site':
  svn:auto-props
    *.html = svn:eol-style=native
\end{verbatim}

Properties inherited from above the root of the working copy are cached
in the working copy\textquotesingle s administrative database when the
working copy is initially checked out and then refreshed whenever the
working copy is updated. This means that you don\textquotesingle t need
access to your repository to view inherited properties. This allows
Subversion subcommands that have traditionally not required access to
the repository (e.g. \texttt{svn\ add} ) to remain "disconnected" while
still accessing properties inherited from paths not found in the working
copy. However it also means that inherited properties from above the
root of the working copy may have changed since your most recent update,
causing your local cache to become out of date. So if you require the
absolute latest value of some inherited property, it\textquotesingle s
always safest to update your working copy first or query the repository
directly.

At this point you might be thinking, "nice trick, but what good is it?"
By itself property inheritance isn\textquotesingle t very useful. Before
1.8, all of Subversion\textquotesingle s own reserved \texttt{svn:*}
properties (and likely all of your own custom user properties) applied
only to the path on which they were set or at most, the
path\textquotesingle s immediate children\footnote{The one noteable
  exception to this being the \texttt{svn:mergeinfo} property, which is
  inheritable---see
  \hyperref[svn.branchmerge.basicmerging.mergeinfo]{Mergeinfo and
  Previews}} Rather, inheritable properties are a tool that Subversion
uses to do other more interesting things, like setting automatic
properties with the \texttt{svn:auto-props} property or repository-wide
ignores with the \texttt{svn:global-ignores} property---see
\hyperref[svn.advanced.props.auto]{Automatic Property Setting} and
\hyperref[svn.advanced.props.special.ignore]{Ignoring Unversioned Items}
for more information about these special properties and how to use them.

Currently inheritable properties are primarily useful only as regards
the \texttt{svn:auto-props} and \texttt{svn:global-ignores} properties
but that doesn\textquotesingle t mean those two properties are the end
of the story. Look for more features to be built with inherited
properties in future releases of Subversion---a log message templating
mechanism comes to mind. In the meantime feel free to use the feature
however you\textquotesingle d like. Any piece of versioned metadata you
want to apply to your whole repository (or large subsections thereof)
can easily be stored in a property on the root of your repository (or
the appropriate subtree). We suspect that some users and administrators
will come up with clever ways to use inheritable properties which we
never considered.

\subsection{Automatic Property Setting}\label{svn.advanced.props.auto}

Properties are a powerful feature of Subversion, acting as key
components of many Subversion features discussed elsewhere in this and
other chapters---textual diff and merge support, keyword substitution,
newline translation, and so on. But to get the full benefit of
properties, they must be set on the right files and directories.
Unfortunately, that step can be easily forgotten in the routine of
things, especially since failing to set a property
doesn\textquotesingle t usually result in an obvious error (at least
compared to, say, failing to add a file to version control). To help
your properties get applied to the places that need them, Subversion
provides a few simple but useful features.

Whenever you introduce a file to version control using the
\texttt{svn\ add} or \texttt{svn\ import} commands, Subversion tries to
assist by setting some common file properties automatically. First, on
operating systems whose filesystems support an execute permission bit,
Subversion will automatically set the \texttt{svn:executable} property
on newly added or imported files whose execute bit is enabled. (See
\hyperref[svn.advanced.props.special.executable]{File Executability}
later in this chapter for more about this property.)

Second, Subversion tries to determine the file\textquotesingle s MIME
type. If you\textquotesingle ve configured a \texttt{mime-types-files}
runtime configuration parameter, Subversion will try to find a MIME type
mapping in that file for your file\textquotesingle s extension. If it
finds such a mapping, it will set your file\textquotesingle s
\texttt{svn:mime-type} property to the MIME type it found. If no mapping
file is configured, or no mapping for your file\textquotesingle s
extension could be found, Subversion will fall back to heuristic
algorithms to determine the file\textquotesingle s MIME type. Depending
on how it is built, Subversion 1.7 can make use of file scanning
libraries\footnote{Currently, libmagic is the support library used to
  accomplish this.} to detect a file\textquotesingle s type based on its
content. Failing all else, Subversion will employ its own very basic
heuristic to determine whether the file contains nontextual content. If
so, it automatically sets the \texttt{svn:mime-type} property on that
file to \texttt{application/octet-stream} (the generic ``this is a
collection of bytes'' MIME type). Of course, if Subversion guesses
incorrectly, or if you wish to set the \texttt{svn:mime-type} property
to something more precise---perhaps \texttt{image/png} or
\texttt{application/x-shockwave-flash}---you can always remove or edit
that property. (For more on Subversion\textquotesingle s use of MIME
types, see \hyperref[svn.advanced.props.special.mime-type]{File Content
Type} later in this chapter.)

UTF-16 is commonly used to encode files whose semantic content is
textual in nature, but the encoding itself makes heavy use of bytes
which are outside the typical ASCII character byte range. As such,
Subversion will tend to classify such files as binary files, much to the
chagrin of users who desire line-based differencing and merging, keyword
substitution, and other behaviors for those files.

Subversion also provides, via its runtime configuration system (see
\hyperref[svn.advanced.confarea]{Runtime Configuration Area}), a more
flexible automatic property setting feature that allows you to create
mappings of filename patterns to property names and values. Once again,
these mappings affect adds and imports, and can not only override the
default MIME type decision made by Subversion during those operations,
but can also set additional Subversion or custom properties, too. For
example, you might create a mapping that says that anytime you add JPEG
files---ones whose names match the pattern \texttt{*.jpg}---Subversion
should automatically set the \texttt{svn:mime-type} property on those
files to \texttt{image/jpeg}. Or perhaps any files that match
\texttt{*.cpp} should have \texttt{svn:eol-style} set to
\texttt{native}, and \texttt{svn:keywords} set to \texttt{Id}. For more
details on automatic property support in the runtime configuration see
\hyperref[svn.advanced.confarea.opts.config]{General configuration}.

While automatic property support via the runtime configuration system is
certainly handy, Subversion administrators might prefer a set of
property definitions which all connecting clients automatically consider
when operating on working copies checked out from a given server.
Subversion 1.8 and newer clients support such functionality through the
\texttt{svn:auto-props} inheritable property.

The \texttt{svn:auto-props} property works like the runtime
configuration to automatically set properties on files when they are
added or imported. The value of the \texttt{svn:auto-props} property is
expected to be the same as the \texttt{auto-props} runtime configuration
option (i.e. Any number of key-value pairs in the format FILE\_PATTERN =
PROPNAME=VALUE{[};PROPNAME=VALUE ...{]}) Like the \texttt{auto-props}
runtime option, the \texttt{svn:auto-props} property can be disregarded
when using the \texttt{-\/-no-auto-props} option, but unlike the config
option, the \texttt{svn:auto-props} property is \emph{not} disabled when
the \texttt{enable-auto-props} configuration option is set to
\texttt{no}.

For example, say you have checked out a working copy of your
\texttt{trunk} branch and need to add a new file (let\textquotesingle s
assume that automatic properties in your runtime configuration are
disabled):

\begin{verbatim}
$ svn st
?       calc/data.c

$ svn add calc/data.c
A         calc/data.c

$ svn proplist -v calc/data.c
Properties on 'calc/data.c':
  svn:eol-style
    native
\end{verbatim}

Notice that after you place the unversioned file \texttt{data.c} under
version control the \texttt{svn:eol-style} property was automatically
set on it. Since we assumed that the \texttt{auto-props} runtime
configuration option is disabled, we know that the
\texttt{svn:auto-props} property must be set on some parent path of
\texttt{data.c}. Using the \texttt{svn\ propget} subcommand with the
\texttt{-\/-show-inherited-props} option we see that this is indeed the
case:

\begin{verbatim}
$ svn propget svn:auto-props --show-inherited-props -v calc
Inherited properties on 'calc',
from 'http://svn.example.com/repos':
  svn:auto-props
    *.py = svn:eol-style=native
    *.c = svn:eol-style=native
    *.h = svn:eol-style=native
\end{verbatim}

Unlike the \texttt{svn:global-ignores} property and its analogous
runtime configuration \texttt{global-ignores}, which are combined, the
\texttt{svn:auto-props} property \emph{overrides} the
\texttt{auto-props} runtime configuration if it defines an auto-prop for
the \emph{same} pattern as the runtime configuration. Automatic
properties inherited\footnote{Remember that users can only inherit
  properties from paths for which they have read access. So if an
  administrator sets \texttt{svn:auto-props} on some high-level parent
  path (e.g. the repository root), they need to be sure all users have
  read access to that path or the desired automatic property setting
  won\textquotesingle t kick in.} from one path can also override the
\emph{identical} pattern inherited from a different path. The hierarchy
of these overrides works as follows:

\begin{itemize}
\item
  An auto-prop, for a given pattern, defined in \texttt{svn:auto-props}
  overrides the same auto-prop for the identical pattern in the
  \texttt{auto-props} runtime configuration.
\item
  If an auto-prop, for a given pattern, is inherited from more than one
  parents\textquotesingle{} \texttt{svn:auto-props} property, the nearer
  path-wise parent overrides the more distant parents.
\item
  An auto-prop, for a given pattern, defined in a
  \texttt{svn:auto-props} property explicitly set on a path overrides
  the same auto-prop(s) for the identical pattern inherited from any
  parents.
\end{itemize}

Let\textquotesingle s look at an example. Suppose you have this runtime
configuration:

\begin{verbatim}
[miscellany]
enable-auto-props = yes
[auto-props]
*.py  = svn:eol-style=CR
*.c   = svn:eol-style=CR
*.h   = svn:eol-style=CR
*.cpp = svn:eol-style=CR
\end{verbatim}

And you want to add three files in the \texttt{calc} directory of your
working copy:

\begin{verbatim}
$ svn st
?       calc/data-binding.cpp
?       calc/data.c
?       calc/editor.py
\end{verbatim}

Let\textquotesingle s check what \texttt{svn:auto-props} apply to
\texttt{calc}:

\begin{verbatim}
$ svn propget svn:auto-props -v --show-inherited-props calc
Inherited properties on 'calc',
from 'http://svn.example.com/repos':
  svn:auto-props
    *.py = svn:eol-style=native
    *.c = svn:eol-style=native
    *.h = svn:eol-style=native

Inherited properties on 'calc',
from '.':
  svn:auto-props
    *.py = svn:eol-style=native
    *.c = svn:keywords=Author Date Id Rev URL
\end{verbatim}

When we add these three files what auto-props do we expect? We add the
trio to version control and then check:

\begin{verbatim}
$ svn add calc --force
A         calc/data-binding.cpp
A         calc/data.c
A         calc/editor.py
\end{verbatim}

The file \texttt{data-binding.cpp} has only one matching pattern,
\texttt{*.cpp\ =\ svn:eol-style=CR} in the runtime configuration, so
obviously the \texttt{svn:eol-style} property is set to \texttt{CR}:

\begin{verbatim}
$ svn proplist -v calc/data-binding.cpp
Properties on 'calc/data-binding.cpp':
  svn:eol-style
    CR
\end{verbatim}

The file \texttt{editor.py} matches a single pattern in runtime config
and both of the \texttt{svn:auto-props} properties, but by the hierarchy
described above, the property explicitly set on \texttt{calc},
\texttt{*.py\ =\ svn:eol-style=native}, takes precedence. So the
\texttt{svn:eol-style} property is set to \texttt{native:}:

\begin{verbatim}
$ svn proplist -v calc/editor.py
Properties on 'calc/editor.py':
  svn:eol-style
    native
\end{verbatim}

The file \texttt{data.c} also matches patterns in the runtime config and
both of the inherited \texttt{svn:auto-props} properties. The
\texttt{svn:keywords} auto-prop is only defined once, on \texttt{calc},
so \texttt{data.c} automatically gets that property. The
\texttt{svn:auto-props} on \texttt{calc} don\textquotesingle t define a
\texttt{svn:eol-style} value however, so the nearest inherited parent,
\texttt{http://svn.example.com/repos}, provides that value:

\begin{verbatim}
$ svn proplist -v calc/data.c
Properties on 'calc/data.c':
  svn:eol-style
    native
  svn:keywords
    Author Date Id Rev URL
\end{verbatim}

Overriding auto-props only applies for \emph{identical} patterns. If a
file to be added or imported matches more than one pattern, then there
is no guarantee which pattern\textquotesingle s auto-props will be
applied. For example, say you want to add the file \texttt{foo.cpp} in
the directory \texttt{bar}. Further, suppose the \texttt{svn:auto-props}
property is set on \texttt{bar} with the value:

\begin{verbatim}
*.c*  = svn:eol-style=native
*.cpp = svn:eol-style=native;svn:keywords=Author Date Id Rev URL
\end{verbatim}

Since \texttt{foo.cpp} matches both patterns, there is no way to know if
the \texttt{svn:keywords} property will be set on \texttt{foo.cpp} when
it is added.

A final note on \texttt{svn:auto-props}. This property (along with the
similar \texttt{svn:global-ignores}, see
\hyperref[svn.advanced.props.special.ignore]{Ignoring Unversioned
Items}) only provides a \emph{recommendation} to clients that understand
the meaning of the property. Older clients will ignore these properties,
the \texttt{-\/-no-auto-props} option will disregard them, a user might
manually change or remove automatic properties after they have been
set---there are numerous ways in which the recommended properties
contained in \texttt{svn:auto-props} can be by-passed. Given this,
administrators will still need to use hook scripts to validate that the
properties added to and modified on files and directories match the
administrator\textquotesingle s preferred policies, rejecting commits
which are non-compliant in this fashion. (See
\hyperref[svn.reposadmin.hooks]{Implementing Repository Hooks} for more
about hook scripts.)

\subsection{Subversion\textquotesingle s Reserved
Properties}\label{svn.advanced.props.ref}

In this section, we\textquotesingle ll briefly summarize all the
properties which Subversion reserves for its own use.
We\textquotesingle ll look at both types of properties---those which are
associated with individual versioned files and directories, and those
which are associated with revisions.

\subsubsection{Versioned
properties}\label{svn.advanced.props.ref.versioned}

These are the versioned (or node) properties that Subversion reserves
for its own use:

\begin{description}
\item[\texttt{svn:auto-props}]
If present on a directory, the value is a set of automatic property
definitions which apply to all files under the directory, See
\hyperref[svn.advanced.props.auto]{Automatic Property Setting}.
\item[\texttt{svn:executable}]
If present on a file, the client will make the file executable in
Unix-hosted working copies. See
\hyperref[svn.advanced.props.special.executable]{File Executability}.
\item[\texttt{svn:mime-type}]
If present on a file, the value indicates the file\textquotesingle s
MIME type. This allows the client to decide whether line-based
contextual merging is safe to perform during updates, and can also
affect how the file behaves when fetched via a web browser. See
\hyperref[svn.advanced.props.special.mime-type]{File Content Type}.
\item[\texttt{svn:ignore}]
If present on a directory, the value is a list of \emph{unversioned}
file patterns to be ignored by \texttt{svn\ status} and other
subcommands. See \hyperref[svn.advanced.props.special.ignore]{Ignoring
Unversioned Items}.
\item[\texttt{svn:global-ignores}]
If present on a directory, the value is a list of \emph{unversioned}
file patterns to be ignored by \texttt{svn\ status} and other
subcommands. Unlike \texttt{svn:ignore} these patterns apply to
\emph{all} unversioned subtrees under the directory, not just the
directory\textquotesingle s immediate file children. See
\hyperref[svn.advanced.props.special.ignore]{Ignoring Unversioned
Items}.
\item[\texttt{svn:keywords}]
If present on a file, the value tells the client how to expand
particular keywords within the file. See
\hyperref[svn.advanced.props.special.keywords]{Keyword Substitution}.
\item[\texttt{svn:eol-style}]
If present on a file, the value tells the client how to manipulate the
file\textquotesingle s line-endings in the working copy and in exported
trees. See \hyperref[svn.advanced.props.special.eol-style]{End-of-Line
Character Sequences} and \hyperref[svn.ref.svn.c.export]{svn export}.
\item[\texttt{svn:externals}]
If present on a directory, the value is a multiline list of other paths
and URLs the client should check out. See
\hyperref[svn.advanced.externals]{Externals Definitions}.
\item[\texttt{svn:special}]
If present on a file, indicates that the file is not an ordinary file,
but a symbolic link or other special object.\footnote{As of this
  writing, symbolic links are indeed the only ``special'' objects. But
  there might be more in future releases of Subversion.}
\item[\texttt{svn:needs-lock}]
If present on a file, tells the client to make the file read-only in the
working copy, as a reminder that the file should be locked before
editing begins. See
\hyperref[svn.advanced.locking.lock-communication]{Lock Communication}.
\item[\texttt{svn:mergeinfo}]
Used by Subversion to track merge data. See
\hyperref[svn.branchmerge.basicmerging.mergeinfo]{Mergeinfo and
Previews} for details, but you should never edit this property unless
you \emph{really} know what you\textquotesingle re doing.
\end{description}

\subsubsection{Unversioned
properties}\label{svn.advanced.props.ref.unversioned}

The following are the unversioned (or revision) properties that
Subversion reserves for its own use. Most of these appear on every
revision in the repository, carrying important information about the
origin and nature of the changes made in that revision.

\begin{description}
\item[\texttt{svn:author}]
If present, contains the authenticated username of the person who
created the revision. (If not present, the revision was committed
anonymously.)
\item[\texttt{svn:autoversioned}]
If present, the revision was created via the autoversioning feature. See
\hyperref[svn.webdav.autoversioning]{Autoversioning}.
\item[\texttt{svn:date}]
Contains the UTC time the revision was created, in ISO 8601 format. The
value comes from the \emph{server} machine\textquotesingle s clock, not
the client\textquotesingle s.
\item[\texttt{svn:log}]
Contains the log message describing the revision.
\end{description}

Certain auxiliary tools in the Subversion toolchain---namely,
\texttt{svnrdump} and \texttt{svnsync}---also use unversioned properties
for their own accounting purposes. These properties are found only on
revision 0 of repositories on which these tools are operating. For more
about \texttt{svnrdump} and \texttt{svnsync} and the functionality they
offer, see \hyperref[svn.reposadmin]{Repository Administration}. The
following are the properties created and managed by these tools.

\begin{description}
\item[\texttt{svn:rdump-lock}]
Used to temporarily enforce mutually exclusive access to the repository
by \texttt{svnrdump\ load}. This property is generally only observed
when such an operation is active---or when an \texttt{svnrdump} command
failed to cleanly disconnect from the repository. (This property is only
relevant when it appears on revision 0.)
\item[\texttt{svn:sync-currently-copying}]
Contains the revision number from the source repository which is
currently being mirrored to this one by the \texttt{svnsync} tool. (This
property is only relevant when it appears on revision 0.)
\item[\texttt{svn:sync-from-uuid}]
Contains the UUID of the repository of which this repository has been
initialized as a mirror by the \texttt{svnsync} tool. (This property is
only relevant when it appears on revision 0.)
\item[\texttt{svn:sync-from-url}]
Contains the URL of the repository directory of which this repository
has been initialized as a mirror by the \texttt{svnsync} tool. (This
property is only relevant when it appears on revision 0.)
\item[\texttt{svn:sync-last-merged-rev}]
Contains the revision of the source repository which was most recently
and successfully mirrored to this one. (This property is only relevant
when it appears on revision 0.)
\item[\texttt{svn:sync-lock}]
Used to temporarily enforce mutually exclusive access to the repository
by \texttt{svnsync} mirroring operations. This property is generally
only observed when such an operation is active---or when an
\texttt{svnsync} command failed to cleanly disconnect from the
repository. (This property is only relevant when it appears on revision
0.)
\end{description}

\section{File Portability}\label{svn.advanced.props.file-portability}

Fortunately for Subversion users who routinely find themselves on
different computers with different operating systems,
Subversion\textquotesingle s command-line program behaves almost
identically on all those systems. If you know how to wield \texttt{svn}
on one platform, you know how to wield it everywhere.

However, the same is not always true of other general classes of
software or of the actual files you keep in Subversion. For example, on
a Windows machine, the definition of a ``text file'' would be similar to
that used on a Linux box, but with a key difference---the character
sequences used to mark the ends of the lines of those files. There are
other differences, too. Unix platforms have (and Subversion supports)
symbolic links; Windows does not. Unix platforms use filesystem
permission to determine executability; Windows uses filename extensions.

Because Subversion is in no position to unite the whole world in common
definitions and implementations of all of these things, the best it can
do is to try to help make your life simpler when you need to work with
your versioned files and directories on multiple computers and operating
systems. This section describes some of the ways Subversion does this.

\subsection{File Content
Type}\label{svn.advanced.props.special.mime-type}

Subversion joins the ranks of the many applications that recognize and
make use of Multipurpose Internet Mail Extensions (MIME) content types.
Besides being a general-purpose storage location for a
file\textquotesingle s content type, the value of the
\texttt{svn:mime-type} file property determines some behavioral
characteristics of Subversion itself.

Identifying File Types

Various programs on most modern operating systems make assumptions about
the type and format of the contents of a file by the
file\textquotesingle s name, specifically its file extension. For
example, files whose names end in \texttt{.txt} are generally assumed to
be human-readable; that is, able to be understood by simple perusal
rather than requiring complex processing to decipher. Files whose names
end in \texttt{.png}, on the other hand, are assumed to be of the
Portable Network Graphics type---not human-readable at all, and sensible
only when interpreted by software that understands the PNG format and
can render the information in that format as a raster image.

Unfortunately, some of those extensions have changed their meanings over
time. When personal computers first appeared, a file named
\texttt{README.DOC} would have almost certainly been a plain-text file,
just like today\textquotesingle s \texttt{.txt} files. But by the
mid-1990s, you could almost bet that a file of that name would not be a
plain-text file at all, but instead a Microsoft Word document in a
proprietary, non-human-readable format. But this change
didn\textquotesingle t occur overnight---there was certainly a period of
confusion for computer users over what exactly they had in hand when
they saw a \texttt{.DOC} file.\footnote{You think that was rough? During
  that same era, WordPerfect also used \texttt{.DOC} for their
  proprietary file format\textquotesingle s preferred extension!}

The popularity of computer networking cast still more doubt on the
mapping between a file\textquotesingle s name and its content. With
information being served across networks and generated dynamically by
server-side scripts, there was often no real file per se, and therefore
no filename. Web servers, for example, needed some other way to tell
browsers what they were downloading so that the browser could do
something intelligent with that information, whether that was to display
the data using a program registered to handle that datatype or to prompt
the user for where on the client machine to store the downloaded data.

Eventually, a standard emerged for, among other things, describing the
contents of a data stream. In 1996, RFC 2045 was published. It was the
first of five RFCs describing MIME. It describes the concept of media
types and subtypes and recommends a syntax for the representation of
those types. Today, MIME media types---or ``MIME types''---are used
almost universally across email applications, web servers, and other
software as the de facto mechanism for clearing up the file content
confusion.

For example, one of the benefits that Subversion typically provides is
contextual, line-based merging of changes received from the server
during an update into your working file. But for files containing
nontextual data, there is often no concept of a ``line.'' So, for
versioned files whose \texttt{svn:mime-type} property is set to a
nontextual MIME type (generally, something that doesn\textquotesingle t
begin with \texttt{text/}, though there are exceptions), Subversion does
not attempt to perform contextual merges during updates. Instead, any
time you have locally modified a binary working copy file that is also
being updated, your file is left untouched and Subversion creates two
new files. One file has a \texttt{.oldrev} extension and contains the
BASE revision of the file. The other file has a \texttt{.newrev}
extension and contains the contents of the updated revision of the file.
This behavior is really for the protection of the user against failed
attempts at performing contextual merges on files that simply cannot be
contextually merged.

Additionally, since the acts of displaying line-based differences and
line-based change attribution are, rather obviously, dependent on there
being a meaningful definition of ``line'' for a given file, files with
nontextual MIME types will by default trigger errors when used as the
targets of \texttt{svn\ diff} and \texttt{svn\ annotate} operations.
This can be especially frustrating for users with XML files whose
\texttt{svn:mime-type} property is set to something such as
\texttt{application/xml} which is not unambiguously human-readable and
as such is treated as nontextual by Subversion. Fortunately, those
subcommands offer a \texttt{-\/-force} option for forcing Subversion to
attempt the operations in spite of the apparent non-human-readability of
the files.

The \texttt{svn:mime-type} property, when set to a value that does not
indicate textual file contents, can cause some unexpected behaviors with
respect to other properties. For example, since the idea of line endings
(and therefore, line-ending conversion) makes no sense when applied to
nontextual files, Subversion will prevent you from setting the
\texttt{svn:eol-style} property on such files. This is obvious when
attempted on a single file target---\texttt{svn\ propset} will error
out. But it might not be as clear if you perform a recursive property
set, where Subversion will silently skip over files that it deems
unsuitable for a given property.

Subversion provides a number of mechanisms by which to automatically set
the \texttt{svn:mime-type} property on a versioned file. See
\hyperref[svn.advanced.props.auto]{Automatic Property Setting} for
details.

Also, if the \texttt{svn:mime-type} property is set, then the Subversion
Apache module will use its value to populate the \texttt{Content-type:}
HTTP header when responding to GET requests. This gives your web browser
a crucial clue about how to display a file when you use it to peruse
your Subversion repository\textquotesingle s contents.

\subsection{File
Executability}\label{svn.advanced.props.special.executable}

On many operating systems, the ability to execute a file as a command is
governed by the presence of an execute permission bit. This bit usually
defaults to being disabled, and must be explicitly enabled by the user
for each file that needs it. But it would be a monumental hassle to have
to remember exactly which files in a freshly checked-out working copy
were supposed to have their executable bits toggled on, and then to have
to do that toggling. So, Subversion provides the \texttt{svn:executable}
property as a way to specify that the executable bit for the file on
which that property is set should be enabled, and Subversion honors that
request when populating working copies with such files.

This property has no effect on filesystems that have no concept of an
executable permission bit, such as FAT32 and NTFS.\footnote{The Windows
  filesystems use file extensions (such as \texttt{.EXE}, \texttt{.BAT},
  and \texttt{.COM}) to denote executable files.} Also, although it has
no defined values, Subversion will force its value to \texttt{*} when
setting this property. Finally, this property is valid only on files,
not on directories.

\subsection{End-of-Line Character
Sequences}\label{svn.advanced.props.special.eol-style}

Unless otherwise noted using a versioned file\textquotesingle s
\texttt{svn:mime-type} property, Subversion assumes the file contains
human-readable data. Generally speaking, Subversion uses this knowledge
only to determine whether contextual difference reports for that file
are possible. Otherwise, to Subversion, bytes are bytes.

This means that by default, Subversion doesn\textquotesingle t pay any
attention to the type of end-of-line (EOL) markers used in your files.
Unfortunately, different operating systems have different conventions
about which character sequences represent the end of a line of text in a
file. For example, the usual line-ending token used by software on the
Windows platform is a pair of ASCII control characters---a carriage
return (\texttt{CR}) followed by a line feed (\texttt{LF}). Unix
software, however, just uses the \texttt{LF} character to denote the end
of a line.

Not all of the various tools on these operating systems understand files
that contain line endings in a format that differs from the native
line-ending style of the operating system on which they are running. So,
typically, Unix programs treat the \texttt{CR} character present in
Windows files as a regular character (usually rendered as
\texttt{\^{}M}), and Windows programs combine all of the lines of a Unix
file into one giant line because no \texttt{CR} characters are found to
denote the ends of the lines.

This sensitivity to foreign EOL markers can be frustrating for folks who
share a file across different operating systems. For example, consider a
source code file, and developers who edit this file on both Windows and
Unix systems. If all the developers always use tools that preserve the
line-ending style of the file, no problems occur.

But in practice, many common tools either fail to properly read a file
with foreign EOL markers, or convert the file\textquotesingle s line
endings to the native style when the file is saved. If the former is
true for a developer, he has to use an external conversion utility (such
as \texttt{dos2unix} or its companion, \texttt{unix2dos}) to prepare the
file for editing. The latter case requires no extra preparation. But
both cases result in a file that differs from the original quite
literally on every line! Prior to committing his changes, the user has
two choices. Either he can use a conversion utility to restore the
modified file to the same line-ending style that it was in before his
edits were made, or he can simply commit the file---new EOL markers and
all.

The result of scenarios like these include wasted time and unnecessary
modifications to committed files. Wasted time is painful enough. But
when commits change every line in a file, this complicates the job of
determining which of those lines were changed in a nontrivial way. Where
was that bug really fixed? On what line was a syntax error introduced?

The solution to this problem is the \texttt{svn:eol-style} property.
When this property is set to a valid value, Subversion uses it to
determine what special processing to perform on the file so that the
file\textquotesingle s line-ending style isn\textquotesingle t
flip-flopping with every commit that comes from a different operating
system. The valid values are:

\begin{description}
\item[\texttt{native}]
This causes the file to contain the EOL markers that are native to the
operating system on which Subversion was run. In other words, if a user
on a Windows machine checks out a working copy that contains a file with
an \texttt{svn:eol-style} property set to \texttt{native}, that file
will contain \texttt{CRLF} EOL markers. A Unix user checking out a
working copy that contains the same file will see \texttt{LF} EOL
markers in his copy of the file.

Note that Subversion will actually store the file in the repository
using normalized \texttt{LF} EOL markers regardless of the operating
system. This is basically transparent to the user, though.
\item[\texttt{CRLF}]
This causes the file to contain \texttt{CRLF} sequences for EOL markers,
regardless of the operating system in use.
\item[\texttt{LF}]
This causes the file to contain \texttt{LF} characters for EOL markers,
regardless of the operating system in use.
\item[\texttt{CR}]
This causes the file to contain \texttt{CR} characters for EOL markers,
regardless of the operating system in use. This line-ending style is not
very common.
\end{description}

\section{Ignoring Unversioned
Items}\label{svn.advanced.props.special.ignore}

In any given working copy, there is a good chance that alongside all
those versioned files and directories are other files and directories
that are neither versioned nor intended to be. Text editors litter
directories with backup files. Software compilers generate
intermediate---or even final---files that you typically
wouldn\textquotesingle t bother to version. And users themselves drop
various other files and directories wherever they see fit, often in
version control working copies.

It\textquotesingle s ludicrous to expect Subversion working copies to be
somehow impervious to this kind of clutter and impurity. In fact,
Subversion counts it as a \emph{feature} that its working copies are
just typical directories, just like unversioned trees. But these
not-to-be-versioned files and directories can cause some annoyance for
Subversion users. For example, because the \texttt{svn\ add} and
\texttt{svn\ import} commands act recursively by default and
don\textquotesingle t know which files in a given tree you do and
don\textquotesingle t wish to version, it\textquotesingle s easy to
accidentally add stuff to version control that you
didn\textquotesingle t mean to. And because \texttt{svn\ status}
reports, by default, every item of interest in a working
copy---including unversioned files and directories---its output can get
quite noisy where many of these things exist.

So Subversion provides several ways for telling it which files you would
prefer that it simply disregard. One of the ways involves the use of
Subversion\textquotesingle s runtime configuration system (see
\hyperref[svn.advanced.confarea]{Runtime Configuration Area}), and
therefore applies to all the Subversion operations that make use of that
runtime configuration---generally those performed on a particular
computer or by a particular user of a computer. Two other methods make
use of Subversion\textquotesingle s directory property support and are
more tightly bound to the versioned tree itself, and therefore affects
everyone who has a working copy of that tree. All of these mechanisms
use file patterns (strings of literal and special wildcard characters
used to match against filenames) to decide which files to ignore.

The Subversion runtime configuration system provides an option,
\texttt{global-ignores}, whose value is a whitespace-delimited
collection of file patterns. The Subversion client checks these patterns
against the names of the files that are candidates for addition to
version control, as well as to unversioned files that the
\texttt{svn\ status} command notices. If any file\textquotesingle s name
matches one of the patterns, Subversion will basically act as if the
file didn\textquotesingle t exist at all. This is really useful for the
kinds of files that you almost never want to version, such as editor
backup files such as Emacs\textquotesingle{} \texttt{*\textasciitilde{}}
and \texttt{.*\textasciitilde{}} files.

File Patterns in Subversion

File patterns (also called globs or shell wildcard patterns) are strings
of characters that are intended to be matched against filenames,
typically for the purpose of quickly selecting some subset of similar
files from a larger grouping without having to explicitly name each
file. The patterns contain two types of characters: regular characters,
which are compared explicitly against potential matches, and special
wildcard characters, which are interpreted differently for matching
purposes.

There are different types of file pattern syntaxes, but Subversion uses
the one most commonly found in Unix systems implemented as the
\texttt{fnmatch} system function. It supports the following wildcards,
described here simply for your convenience:

\begin{description}
\item[\texttt{?}]
Matches any single character
\item[\texttt{*}]
Matches any string of characters, including the empty string
\item[\texttt{{[}}]
Begins a character class definition terminated by \texttt{{]}}, used for
matching a subset of characters
\end{description}

You can see this same pattern matching behavior at a Unix shell prompt.
The following are some examples of patterns being used for various
things:

\begin{verbatim}
$ ls   ### the book sources
appa-quickstart.xml             ch06-server-configuration.xml
appb-svn-for-cvs-users.xml      ch07-customizing-svn.xml
appc-webdav.xml                 ch08-embedding-svn.xml
book.xml                        ch09-reference.xml
ch00-preface.xml                ch10-world-peace-thru-svn.xml
ch01-fundamental-concepts.xml   copyright.xml
ch02-basic-usage.xml            foreword.xml
ch03-advanced-topics.xml        images/
ch04-branching-and-merging.xml  index.xml
ch05-repository-admin.xml       styles.css
$ ls ch*   ### the book chapters
ch00-preface.xml                ch06-server-configuration.xml
ch01-fundamental-concepts.xml   ch07-customizing-svn.xml
ch02-basic-usage.xml            ch08-embedding-svn.xml
ch03-advanced-topics.xml        ch09-reference.xml
ch04-branching-and-merging.xml  ch10-world-peace-thru-svn.xml
ch05-repository-admin.xml
$ ls ch?0-*   ### the book chapters whose numbers end in zero
ch00-preface.xml  ch10-world-peace-thru-svn.xml
$ ls ch0[3578]-*   ### the book chapters that Mike is responsible for
ch03-advanced-topics.xml   ch07-customizing-svn.xml
ch05-repository-admin.xml  ch08-embedding-svn.xml
$
\end{verbatim}

File pattern matching is a bit more complex than what
we\textquotesingle ve described here, but this basic usage level tends
to suit the majority of Subversion users.

When found on a versioned directory, the \texttt{svn:ignore} property is
expected to contain a list of newline-delimited file patterns that
Subversion should use to determine ignorable objects in that \emph{same}
directory. These patterns do not override those found in the
\texttt{global-ignores} runtime configuration option, but are instead
appended to that list. And it\textquotesingle s worth noting again that,
unlike the \texttt{global-ignores} option, the patterns found in the
\texttt{svn:ignore} property apply only to the directory on which that
property is set, and not to any of its subdirectories. The
\texttt{svn:ignore} property is a good way to tell Subversion to ignore
files that are likely to be present in every user\textquotesingle s
working copy of that directory, such as compiler output or---to use an
example more appropriate to this book---the HTML, PDF, or PostScript
files generated as the result of a conversion of some source DocBook XML
files to a more legible output format.

Subversion 1.8 provides a more powerful version of the
\texttt{svn:ignore} property, the \texttt{svn:global-ignores} property.
Like the \texttt{svn:ignore} property, \texttt{svn:global-ignores} can
only be set on a directory and contains file patterns Subversion uses to
determine ignorable objects.\footnote{The ignore patterns in the
  \texttt{svn:global-ignores} property may be delimited with any
  whitespace (similar to the \texttt{global-ignores} runtime
  configuration option), not just newlines (as with the
  \texttt{svn:ignore} property).} These ignore patterns are also
appended to any patterns defined in the \texttt{global-ignores} runtime
configuration option together with any \texttt{svn:ignore} defined
patterns. Unlike \texttt{svn:ignore} however, the
\texttt{svn:global-ignores} property is inheritable \footnote{Of course
  only a 1.8 or newer Subversion client will recognize the
  inheritability and special meaning of the \texttt{svn:global-ignores}
  property!} and applies to \emph{all} paths under the directory on
which the property is set, not just the immediate children of the
directory.

Subversion\textquotesingle s support for ignorable file patterns extends
only to the one-time process of adding unversioned files and directories
to version control. Once an object is under Subversion\textquotesingle s
control, the ignore pattern mechanisms no longer apply to it. In other
words, don\textquotesingle t expect Subversion to avoid committing
changes you\textquotesingle ve made to a versioned file simply because
that file\textquotesingle s name matches an ignore pattern---Subversion
\emph{always} notices all of its versioned objects.

Ignore Patterns for CVS Users

The Subversion \texttt{svn:ignore} property is very similar in syntax
and function to the CVS \texttt{.cvsignore} file. In fact, if you are
migrating a CVS working copy to Subversion, you can directly migrate the
ignore patterns by using the \texttt{.cvsignore} file as input to the
\texttt{svn\ propset} command:

\begin{verbatim}
$ svn propset svn:ignore -F .cvsignore .
property 'svn:ignore' set on '.'
$
\end{verbatim}

There are, however, some differences in the ways that CVS and Subversion
handle ignore patterns. The two systems use the ignore patterns at some
different times, and there are slight discrepancies in what the ignore
patterns apply to. Also, Subversion does not recognize the use of the
\texttt{!} pattern as a reset back to having no ignore patterns at all.

The ignore patterns in the \texttt{global-ignores} runtime configuration
option tend to be more a matter of personal taste\footnote{Despite being
  a matter of personal taste, if you don\textquotesingle t explicitly
  set the \texttt{global-ignores} runtime configuration option---either
  to your preferred set of patterns or to an empty string---Subversion
  uses a default value. See the \texttt{global-ignores} entry in
  \hyperref[svn.advanced.confarea.opts.config]{General configuration}}
and ties more closely to a user\textquotesingle s particular tool chain
than to the details of any particular working copy\textquotesingle s
needs. So, the rest of this section will focus on the
\texttt{svn:ignore} and \texttt{svn:global-ignores} properies and their
uses.

Say you have the following output from \texttt{svn\ status}:

\begin{verbatim}
$ svn status calc
 M      calc/button.c
?       calc/calculator
?       calc/data.c
?       calc/debug_log
?       calc/debug_log.1
?       calc/debug_log.2.gz
?       calc/debug_log.3.gz
\end{verbatim}

In this example, you have made some property modifications to
\texttt{button.c}, but in your working copy, you also have some
unversioned files: the latest \texttt{calculator} program that
you\textquotesingle ve compiled from your source code, a source file
named \texttt{data.c}, and a set of debugging output logfiles. Now, you
know that your build system always results in the \texttt{calculator}
program being generated.\footnote{Isn\textquotesingle t that the whole
  point of a build system?} And you know that your test suite always
leaves those debugging logfiles lying around. These facts are true for
all working copies of this project, not just your own. And you know that
you aren\textquotesingle t interested in seeing those things every time
you run \texttt{svn\ status}, and you are pretty sure that nobody else
is interested in them either. So you use
\texttt{svn\ propedit\ svn:ignore\ calc} to add some ignore patterns to
the \texttt{calc} directory.

\begin{verbatim}
$ svn propget svn:ignore calc
calculator
debug_log*
$
\end{verbatim}

After you\textquotesingle ve added this property, you will now have a
local property modification on the \texttt{calc} directory. But notice
what else is different about your \texttt{svn\ status} output:

\begin{verbatim}
$ svn status
 M      calc
 M      calc/button.c
?       calc/data.c
\end{verbatim}

Now, all that cruft is missing from the output! Your \texttt{calculator}
compiled program and all those logfiles are still in your working copy;
Subversion just isn\textquotesingle t constantly reminding you that they
are present and unversioned. And now with all the uninteresting noise
removed from the display, you are left with more intriguing items---such
as that source code file \texttt{data.c} that you probably forgot to add
to version control.

Of course, this less-verbose report of your working copy status
isn\textquotesingle t the only one available. If you actually want to
see the ignored files as part of the status report, you can pass the
\texttt{-\/-no-ignore} option to Subversion:

\begin{verbatim}
$ svn status --no-ignore
 M      calc
 M      calc/button.c
I       calc/calculator
?       calc/data.c
I       calc/debug_log
I       calc/debug_log.1
I       calc/debug_log.2.gz
I       calc/debug_log.3.gz
I       calc/wip.1.diff
\end{verbatim}

All of your previously hidden unversioned paths are once again shown,
but now with the \texttt{\textquotesingle{}I\textquotesingle{}\ Ignored}
status. But wait, what about \texttt{wip.1.diff}? The
\texttt{svn:ignore} property on \texttt{calc} doesn\textquotesingle t
include any pattern that matches that filename, so why is it
ignored?\footnote{Let\textquotesingle s assume that you
  don\textquotesingle t have a matching pattern anywhere in your
  \texttt{global-ignores} runtime configuration.} The answer lies in the
third method by which Subversion can disregard unversioned paths, the
inheritable \texttt{svn:global-ignores} property. Using the
\texttt{svn\ propget} subcommand with the
\texttt{-\/-show-inherited-props} option, you see that the
\texttt{svn:global-ignores} property is set on the root of your working
copy, and sure enough, it defines a matching ignore pattern:

\begin{verbatim}
$ svn pg svn:global-ignores calc -v --show-inherited-props
Inherited properties on 'calc',
from '.':
  svn:global-ignores
    *.diff
    *.patch
\end{verbatim}

As mentioned earlier, the list of file patterns to ignore is also used
by \texttt{svn\ add} and \texttt{svn\ import}. Both of these operations
involve asking Subversion to begin managing some set of files and
directories. Rather than force the user to pick and choose which files
in a tree she wishes to start versioning, Subversion uses the ignore
patterns---the global, per-directory, and inherited lists---to determine
which files should not be swept into the version control system as part
of a larger recursive addition or import operation. And here again, you
can use the \texttt{-\/-no-ignore} option to tell Subversion to
disregard its ignores list and operate on all the files and directories
present.

Even if \texttt{svn:ignore} or \texttt{svn:global-ignores} is set, you
may run into problems if you use shell wildcards in a command. Shell
wildcards are expanded into an explicit list of targets before
Subversion operates on them, so running \texttt{svn\ SUBCOMMAND\ *} is
just like running \texttt{svn\ SUBCOMMAND\ file1\ file2\ file3\ …}. In
the case of the \texttt{svn\ add} command, this has an effect similar to
passing the \texttt{-\/-no-ignore} option. So instead of using a
wildcard, use \texttt{svn\ add\ -\/-force\ .} to do a bulk scheduling of
unversioned things for addition. The explicit target will ensure that
the current directory isn\textquotesingle t overlooked because of being
already under version control, and the \texttt{-\/-force} option will
cause Subversion to crawl through that directory, adding unversioned
files while still honoring the \texttt{svn:ignore} and
\texttt{svn:global-ignores} properties and the \texttt{global-ignores}
runtime configuration variable. Be sure to also provide the
\texttt{-\/-depth\ files} option to the \texttt{svn\ add} command if you
don\textquotesingle t want a fully recursive crawl for things to add.

\section{Keyword
Substitution}\label{svn.advanced.props.special.keywords}

Subversion has the ability to substitute keywords---pieces of useful,
dynamic information about a versioned file---into the contents of the
file itself. Keywords generally provide information about the last
modification made to the file. Because this information changes each
time the file changes, and more importantly, just \emph{after} the file
changes, it is a hassle for any process except the version control
system to keep the data completely up to date. Left to human authors,
the information would inevitably grow stale.

For example, say you have a document in which you would like to display
the last date on which it was modified. You could burden every author of
that document to, just before committing their changes, also tweak the
part of the document that describes when it was last changed. But sooner
or later, someone would forget to do that. Instead, simply ask
Subversion to perform keyword substitution on the
\texttt{LastChangedDate} keyword. You control where the keyword is
inserted into your document by placing a keyword anchor at the desired
location in the file. This anchor is just a string of text formatted as
\texttt{\$}\textless KeywordName\textgreater{}\texttt{\$}.

Adding keyword anchor text alone to your file does nothing special.
Subversion will never attempt to perform textual substitutions on your
file contents unless explicitly asked to do so. After all, you might be
writing a document\footnote{\ldots{} or maybe even a section of a book
  \ldots{}} about how to use keywords, and you don\textquotesingle t
want Subversion to substitute your beautiful examples of unsubstituted
keyword anchors!

To tell Subversion whether to substitute keywords on a particular file,
we again turn to the property-related subcommands. The
\texttt{svn:keywords} property, when set on a versioned file, controls
which keywords will be substituted on that file. The value is a
space-delimited list of keyword names or aliases.

For example, say you have a versioned file named \texttt{weather.txt}
that looks like this:

\begin{verbatim}
Here is the latest report from the front lines.
$LastChangedDate$
$Rev$
Cumulus clouds are appearing more frequently as summer approaches.
\end{verbatim}

With no \texttt{svn:keywords} property set on that file, Subversion will
do nothing special. Now, let\textquotesingle s enable substitution of
the \texttt{LastChangedDate} keyword.

\begin{verbatim}
$ svn propset svn:keywords "Date Author" weather.txt
property 'svn:keywords' set on 'weather.txt'
$
\end{verbatim}

Now you have made a local property modification on the
\texttt{weather.txt} file. You will see no changes to the
file\textquotesingle s contents (unless you made some of your own prior
to setting the property). Notice that the file contained a keyword
anchor for the \texttt{Rev} keyword, yet we did not include that keyword
in the property value we set. Subversion will happily ignore requests to
substitute keywords that are not present in the file and will not
substitute keywords that are not present in the \texttt{svn:keywords}
property value.

Immediately after you commit this property change, Subversion will
update your working file with the new substitute text. Instead of seeing
your keyword anchor \texttt{\$LastChangedDate\$}, you\textquotesingle ll
see its substituted result. That result also contains the name of the
keyword and continues to be delimited by the dollar sign (\texttt{\$})
characters. And as we predicted, the \texttt{Rev} keyword was not
substituted because we didn\textquotesingle t ask for it to be.

Note also that we set the \texttt{svn:keywords} property to
\texttt{Date\ Author}, yet the keyword anchor used the alias
\texttt{\$LastChangedDate\$} and still expanded correctly:

\begin{verbatim}
Here is the latest report from the front lines.
$LastChangedDate: 2006-07-22 21:42:37 -0700 (Sat, 22 Jul 2006) $
$Rev$
Cumulus clouds are appearing more frequently as summer approaches.
\end{verbatim}

If someone else now commits a change to \texttt{weather.txt}, your copy
of that file will continue to display the same substituted keyword value
as before---until you update your working copy. At that time, the
keywords in your \texttt{weather.txt} file will be resubstituted with
information that reflects the most recent known commit to that file.

All keywords are case-sensitive where they appear as anchors in files:
you must use the correct capitalization for the keyword to be expanded.
You should consider the value of the \texttt{svn:keywords} property to
be case-sensitive, too---for the sake of backward compatibility, certain
keyword names will be recognized regardless of case, but this behavior
is deprecated.

Subversion defines the list of keywords available for substitution. That
list contains the following keywords, some of which have aliases that
you can also use:

\begin{description}
\item[\texttt{Date}]
This keyword describes the last time the file was known to have been
changed in the repository, and is of the form
\texttt{\$Date:\ 2006-07-22\ 21:42:37\ -0700\ (Sat,\ 22\ Jul\ 2006)\ \$}.
It may also be specified as \texttt{LastChangedDate}. Unlike the
\texttt{Id} keyword, which uses UTC, the \texttt{Date} keyword displays
dates using the local time zone.
\item[\texttt{Revision}]
This keyword describes the last known revision in which this file
changed in the repository, and looks something like
\texttt{\$Revision:\ 144\ \$}. It may also be specified as
\texttt{LastChangedRevision} or \texttt{Rev}.
\item[\texttt{Author}]
This keyword describes the last known user to change this file in the
repository, and looks something like \texttt{\$Author:\ harry\ \$}. It
may also be specified as \texttt{LastChangedBy}.
\item[\texttt{HeadURL}]
This keyword describes the full URL to the latest version of the file in
the repository, and looks something like
\texttt{\$HeadURL:\ http://svn.example.com/repos/trunk/calc.c\ \$}. It
may be abbreviated as \texttt{URL}.
\item[\texttt{Id}]
This keyword is a compressed combination of the other keywords. Its
substitution looks something like
\texttt{\$Id:\ calc.c\ 148\ 2006-07-28\ 21:30:43Z\ sally\ \$}, and is
interpreted to mean that the file \texttt{calc.c} was last changed in
revision 148 on the evening of July 28, 2006 by the user \texttt{sally}.
The date displayed by this keyword is in UTC, unlike that of the
\texttt{Date} keyword (which uses the local time zone).
\item[\texttt{Header}]
This keyword is similar to the \texttt{Id} keyword but contains the full
URL of the latest revision of the item, identical to \texttt{HeadURL}.
Its substitution looks something like
\texttt{\$Header:\ http://svn.example.com/repos/trunk/calc.c\ 148\ 2006-07-28\ 21:30:43Z\ sally\ \$}.
\end{description}

Several of the preceding descriptions use the phrase ``last known'' or
similar wording. Keep in mind that keyword expansion is a client-side
operation, and your client ``knows'' only about changes that have
occurred in the repository when you update your working copy to include
those changes. If you never update your working copy, your keywords will
never expand to different values even if those versioned files are being
changed regularly in the repository.

Where\textquotesingle s \$GlobalRev\$?

New users are often confused by how the \texttt{\$Rev\$} keyword works.
Since the repository has a single, globally increasing revision number,
many people assume that it is this number that is reflected by the
\texttt{\$Rev\$} keyword\textquotesingle s value. But \texttt{\$Rev\$}
expands to show the last revision in which the file \emph{changed}, not
the last revision to which it was updated. Understanding this clears the
confusion, but frustration often remains---without the support of a
Subversion keyword to do so, how can you automatically get the global
revision number into your files?

To do this, you need external processing. Subversion ships with a tool
called \texttt{svnversion}, which was designed for just this purpose. It
crawls your working copy and generates as output the revision(s) it
finds. You can use this program, plus some additional tooling, to embed
that revision information into your files. For more information on
\texttt{svnversion}, see \hyperref[svn.ref.svnversion]{svnversion
Reference---Subversion Working Copy Version Info}.

In addition to previous set of stock keyword definitions and aliases,
Subversion 1.8 allows you the freedom to define and use custom keywords.
To define a custom keyword, add a token to the value of the
\texttt{svn:keywords} property which is of the form
\texttt{MyKeyword=FORMAT}, where \textless MyKeyword\textgreater{} is
the keyword name (which you\textquotesingle ll use in the keyword
anchor) and \textless FORMAT\textgreater{} is a format string into which
information will be substituted when your keyword is expanded inside
your file.

The format string syntax used for custom keywords supports the following
format codes:

\begin{description}
\item[\texttt{\%a}]
The author of the revision given by \texttt{\%r}.
\item[\texttt{\%b}]
The basename of the URL of the file.
\item[\texttt{\%d}]
Short format of the date of the revision given by \texttt{\%r}.
\item[\texttt{\%D}]
Long format of the date of the revision given by \texttt{\%r}.
\item[\texttt{\%P}]
The file\textquotesingle s path, relative to the repository root.
\item[\texttt{\%r}]
The last known revision in which this file changed in the repository.
(This is the same revision which would be substituted for the
\texttt{Revision} keyword.)
\item[\texttt{\%R}]
The URL to the root of the repository.
\item[\texttt{\%u}]
The URL of the file.
\item[\texttt{\%\_}]
A space character. (Keyword definitions cannot contain a literal space
character.)
\item[\texttt{\%\%}]
A literal percent sign (\textquotesingle{}\texttt{\%}\textquotesingle).
\item[\texttt{\%H}]
Equivalent to \texttt{\%P\%\_\%r\%\_\%d\%\_\%a}.
\item[\texttt{\%I}]
Equivalent to \texttt{\%b\%\_\%r\%\_\%d\%\_\%a}.
\end{description}

As you can see, many of the individual format codes serve as
placeholders for the same information available through the stock
keywords. But of course, the custom keyword format allows you to more
flexibly string together multiple bits of information. For example, you
might wish to have a single keyword in your files which reports the
repository relative path of the file and last-changed revision,
formatted in a pleasant, human-readable way. To do so,
you\textquotesingle d first define your custom keyword:

\begin{verbatim}
$ svn pset svn:keywords "PathRev=%P,%_r%r" calc/button.c
property 'svn:keywords' set on 'button.c'
$
\end{verbatim}

Next, you\textquotesingle d edit the file\textquotesingle s contents to
add the keyword anchor for your custom keyword, which in this case is
\texttt{\$PathRev\$}. After committing these changes, an examination of
your file\textquotesingle s contents will show that your custom keyword
was substituted as you would expect---where previously the file
contained \texttt{\$PathRev\$}, it now reads
\texttt{\$PathRev:\ trunk/calc/button.c,\ r23\ \$}.

Subversion will automatically truncate any keyword expansions which
exceed 255 bytes in length. Also custom keywords defined with names that
exceed 255 bytes will be ignored altogether.

You can also instruct Subversion to maintain a fixed length (in terms of
the number of bytes consumed) for the substituted keyword. By using a
double colon (\texttt{::}) after the keyword name, followed by a number
of space characters, you define that fixed width. When Subversion goes
to substitute your keyword for the keyword and its value, it will
essentially replace only those space characters, leaving the overall
width of the keyword field unchanged. If the substituted value is
shorter than the defined field width, there will be extra padding
characters (spaces) at the end of the substituted field; if it is too
long, it is truncated with a special hash (\texttt{\#}) character just
before the final dollar sign terminator.

For example, say you have a document in which you have some section of
tabular data reflecting the document\textquotesingle s Subversion
keywords. Using the original Subversion keyword substitution syntax,
your file might look something like:

\begin{verbatim}
$Rev$:     Revision of last commit
$Author$:  Author of last commit
$Date$:    Date of last commit
\end{verbatim}

Now, that looks nice and tabular at the start of things. But when you
then commit that file (with keyword substitution enabled, of course),
you see:

\begin{verbatim}
$Rev: 12 $:     Revision of last commit
$Author: harry $:  Author of last commit
$Date: 2006-03-15 02:33:03 -0500 (Wed, 15 Mar 2006) $:    Date of last commit
\end{verbatim}

The result is not so beautiful. And you might be tempted to then adjust
the file after the substitution so that it again looks tabular. But that
holds only as long as the keyword values are the same width. If the last
committed revision rolls into a new place value (say, from 99 to 100),
or if another person with a longer username commits the file, stuff gets
all crooked again. However, if you are using Subversion 1.2 or later,
you can use the new fixed-length keyword syntax and define some field
widths that seem sane, so your file might look like this:

\begin{verbatim}
$Rev::               $:  Revision of last commit
$Author::            $:  Author of last commit
$Date::              $:  Date of last commit
\end{verbatim}

You commit this change to your file. This time, Subversion notices the
new fixed-length keyword syntax and maintains the width of the fields as
defined by the padding you placed between the double colon and the
trailing dollar sign. After substitution, the width of the fields is
completely unchanged---the short values for \texttt{Rev} and
\texttt{Author} are padded with spaces, and the long \texttt{Date} field
is truncated by a hash character:

\begin{verbatim}
$Rev:: 13            $:  Revision of last commit
$Author:: harry      $:  Author of last commit
$Date:: 2006-03-15 0#$:  Date of last commit
\end{verbatim}

The use of fixed-length keywords is especially handy when performing
substitutions into complex file formats that themselves use fixed-length
fields for data, or for which the stored size of a given data field is
overbearingly difficult to modify from outside the
format\textquotesingle s native application. Of course, where binary
file formats are concerned, you must always take great care that any
keyword substitution you introduce---fixed-length or otherwise---does
not violate the integrity of that format. While it might sound easy
enough, this can be an astonishingly difficult task for most of the
popular binary file formats in use today, and \emph{not} something to be
undertaken by the faint of heart!

Be aware that because the width of a keyword field is measured in bytes,
the potential for corruption of multibyte values exists. For example, a
username that contains some multibyte UTF-8 characters might suffer
truncation in the middle of the string of bytes that make up one of
those characters. The result will be a mere truncation when viewed at
the byte level, but will likely appear as a string with an incorrect or
garbled final character when viewed as UTF-8 text. It is conceivable
that certain applications, when asked to load the file, would notice the
broken UTF-8 text and deem the entire file corrupt, refusing to operate
on the file altogether. So, when limiting keywords to a fixed size,
choose a size that allows for this type of byte-wise expansion.

\section{Sparse Directories}\label{svn.advanced.sparsedirs}

By default, most Subversion operations on directories act in a recursive
manner. For example, \texttt{svn\ checkout} creates a working copy with
every file and directory in the specified area of the repository,
descending recursively through the repository tree until the entire
structure is copied to your local disk. Subversion 1.5 introduces a
feature called sparse directories (or shallow checkouts) that allows you
to easily check out a working copy---or a portion of a working
copy---more shallowly than full recursion, with the freedom to bring in
previously ignored files and subdirectories at a later time.

For example, say we have a repository with a tree of files and
directories with names of the members of a human family with pets.
(It\textquotesingle s an odd example, to be sure, but bear with us.) A
regular \texttt{svn\ checkout} operation will give us a working copy of
the whole tree:

\begin{verbatim}
$ svn checkout file:///var/svn/repos mom
A    mom/son
A    mom/son/grandson
A    mom/daughter
A    mom/daughter/granddaughter1
A    mom/daughter/granddaughter1/bunny1.txt
A    mom/daughter/granddaughter1/bunny2.txt
A    mom/daughter/granddaughter2
A    mom/daughter/fishie.txt
A    mom/kitty1.txt
A    mom/doggie1.txt
Checked out revision 1.
$
\end{verbatim}

Now, let\textquotesingle s check out the same tree again, but this time
we\textquotesingle ll ask Subversion to give us only the topmost
directory with none of its children at all:

\begin{verbatim}
$ svn checkout file:///var/svn/repos mom-empty --depth empty
Checked out revision 1
$
\end{verbatim}

Notice that we added to our original \texttt{svn\ checkout} command line
a new \texttt{-\/-depth} option. This option is present on many of
Subversion\textquotesingle s subcommands and is similar to the
\texttt{-\/-non-recursive} (\texttt{-N}) and \texttt{-\/-recursive}
(\texttt{-R}) options. In fact, it combines, improves upon, supercedes,
and ultimately obsoletes these two older options. For starters, it
expands the supported degrees of depth specification available to users,
adding some previously unsupported (or inconsistently supported) depths.
Here are the depth values that you can request for a given Subversion
operation:

\begin{description}
\item[\texttt{-\/-depth\ empty}]
Include only the immediate target of the operation, not any of its file
or directory children.
\item[\texttt{-\/-depth\ files}]
Include the immediate target of the operation and any of its immediate
file children.
\item[\texttt{-\/-depth\ immediates}]
Include the immediate target of the operation and any of its immediate
file or directory children. The directory children will themselves be
empty.
\item[\texttt{-\/-depth\ infinity}]
Include the immediate target, its file and directory children, its
children\textquotesingle s children, and so on to full recursion.
\end{description}

Of course, merely combining two existing options into one hardly
constitutes a new feature worthy of a whole section in our book.
Fortunately, there is more to this story. This idea of depth extends not
just to the operations you perform with your Subversion client, but also
as a description of a working copy citizen\textquotesingle s ambient
depth, which is the depth persistently recorded by the working copy for
that item. Its key strength is this very persistence---the fact that it
is ``sticky''. The working copy remembers the depth
you\textquotesingle ve selected for each item in it until you later
change that depth selection; by default, Subversion commands operate on
the working copy citizens present, regardless of their selected depth
settings.

You can check the recorded ambient depth of a working copy using the
\texttt{svn\ info} command. If the ambient depth is anything other than
infinite recursion, \texttt{svn\ info} will display a line describing
that depth value:

\begin{verbatim}
$ svn info mom-immediates | grep "^Depth:"
Depth: immediates
$
\end{verbatim}

Our previous examples demonstrated checkouts of infinite depth (the
default for \texttt{svn\ checkout}) and empty depth.
Let\textquotesingle s look now at examples of the other depth values:

\begin{verbatim}
$ svn checkout file:///var/svn/repos mom-files --depth files
A    mom-files/kitty1.txt
A    mom-files/doggie1.txt
Checked out revision 1.
$ svn checkout file:///var/svn/repos mom-immediates --depth immediates
A    mom-immediates/son
A    mom-immediates/daughter
A    mom-immediates/kitty1.txt
A    mom-immediates/doggie1.txt
Checked out revision 1.
$
\end{verbatim}

As described, each of these depths is something more than only the
target, but something less than full recursion.

We\textquotesingle ve used \texttt{svn\ checkout} as an example here,
but you\textquotesingle ll find the \texttt{-\/-depth} option present on
many other Subversion commands, too. In those other commands, depth
specification is a way to limit the scope of an operation to some depth,
much like the way the older \texttt{-\/-non-recursive} (\texttt{-N}) and
\texttt{-\/-recursive} (\texttt{-R}) options behave. This means that
when operating on a working copy of some depth, while requesting an
operation of a shallower depth, the operation is limited to that
shallower depth. In fact, we can make an even more general statement:
given a working copy of any arbitrary---even mixed---ambient depth, and
a Subversion command with some requested operational depth, the command
will maintain the ambient depth of the working copy members while still
limiting the scope of the operation to the requested (or default)
operational depth.

In addition to the \texttt{-\/-depth} option, the \texttt{svn\ update}
and \texttt{svn\ switch} subcommands also accept a second depth-related
option: \texttt{-\/-set-depth}. It is with this option that you can
change the sticky depth of a working copy item. Watch what happens as we
take our empty-depth checkout and gradually telescope it deeper using
\texttt{svn\ update\ -\/-set-depth\ NEW-DEPTH\ TARGET}:

\begin{verbatim}
$ svn update --set-depth files mom-empty
Updating 'mom-empty':
A    mom-empty/kittie1.txt
A    mom-empty/doggie1.txt
Updated to revision 1.
$ svn update --set-depth immediates mom-empty
Updating 'mom-empty':
A    mom-empty/son
A    mom-empty/daughter
Updated to revision 1.
$ svn update --set-depth infinity mom-empty
Updating 'mom-empty':
A    mom-empty/son/grandson
A    mom-empty/daughter/granddaughter1
A    mom-empty/daughter/granddaughter1/bunny1.txt
A    mom-empty/daughter/granddaughter1/bunny2.txt
A    mom-empty/daughter/granddaughter2
A    mom-empty/daughter/fishie1.txt
Updated to revision 1.
$
\end{verbatim}

As we gradually increased our depth selection, the repository gave us
more pieces of our tree.

In our example, we operated only on the root of our working copy,
changing its ambient depth value. But we can independently change the
ambient depth value of \emph{any} subdirectory inside the working copy,
too. Careful use of this ability allows us to flesh out only certain
portions of the working copy tree, leaving other portions absent
altogether (hence the ``sparse'' bit of the feature\textquotesingle s
name). Here\textquotesingle s an example of how we might build out a
portion of one branch of our family\textquotesingle s tree, enable full
recursion on another branch, and keep still other pieces pruned (absent
from disk).

\begin{verbatim}
$ rm -rf mom-empty
$ svn checkout file:///var/svn/repos mom-empty --depth empty
Checked out revision 1.
$ svn update --set-depth empty mom-empty/son
Updating 'mom-empty/son':
A    mom-empty/son
Updated to revision 1.
$ svn update --set-depth empty mom-empty/daughter
Updating 'mom-empty/daughter':
A    mom-empty/daughter
Updated to revision 1.
$ svn update --set-depth infinity mom-empty/daughter/granddaughter1
Updating 'mom-empty/daughter/granddaughter1':
A    mom-empty/daughter/granddaughter1
A    mom-empty/daughter/granddaughter1/bunny1.txt
A    mom-empty/daughter/granddaughter1/bunny2.txt
Updated to revision 1.
$
\end{verbatim}

Fortunately, having a complex collection of ambient depths in a single
working copy doesn\textquotesingle t complicate the way you interact
with that working copy. You can still make, revert, display, and commit
local modifications in your working copy without providing any new
options (including \texttt{-\/-depth} and \texttt{-\/-set-depth}) to the
relevant subcommands. Even \texttt{svn\ update} works as it does
elsewhere when no specific depth is provided---it updates the working
copy targets that are present while honoring their sticky depths.

You might at this point be wondering, ``So what? When would I use
this?'' One scenario where this feature finds utility is tied to a
particular repository layout, specifically where you have many related
or codependent projects or software modules living as siblings in a
single repository location (\texttt{trunk/project1},
\texttt{trunk/project2}, \texttt{trunk/project3}, etc.). In such
scenarios, it might be the case that you personally care about only a
handful of those projects---maybe some primary project and a few other
modules on which it depends. You can check out individual working copies
of all of these things, but those working copies are disjoint and, as a
result, it can be cumbersome to perform operations across several or all
of them at the same time. The alternative is to use the sparse
directories feature, building out a single working copy that contains
only the modules you care about. You\textquotesingle d start with an
empty-depth checkout of the common parent directory of the projects, and
then update with infinite depth only the items you wish to have, like we
demonstrated in the previous example. Think of it like an opt-in system
for working copy citizens.

The original (Subversion 1.5) implementation of shallow checkouts was
good, but didn\textquotesingle t support de-telescoping of working copy
items. Subversion 1.6 remedied this problem. For example, running
\texttt{svn\ update\ -\/-set-depth\ empty} in an infinite-depth working
copy will discard everything but the topmost directory.\footnote{Safely,
  of course. As in other situations, Subversion will leave on disk any
  files you\textquotesingle ve modified or which aren\textquotesingle t
  versioned.} Subversion 1.6 also introduced another supported value for
the \texttt{-\/-set-depth} option: \texttt{exclude}. Using
\texttt{-\/-set-depth\ exclude} with \texttt{svn\ update} will cause the
update target to be removed from the working copy entirely---a directory
target won\textquotesingle t even be left present-but-empty. This is
especially handy when there are more things that you\textquotesingle d
like to keep in a working copy than things you\textquotesingle d like to
\emph{not} keep.

Consider a directory with hundreds of subdirectories, one of which you
would like to omit from your working copy. Using an ``additive''
approach to sparse directories, you might check out the directory with
an empty depth, then explicitly telescope (using
\texttt{svn\ update\ -\/-set-depth\ infinity}) each and every
subdirectory of the directory except the one you don\textquotesingle t
care about.

\begin{verbatim}
$ svn checkout http://svn.example.com/repos/many-dirs --depth empty
…
$ svn update --set-depth infinity many-dirs/wanted-dir-1
…
$ svn update --set-depth infinity many-dirs/wanted-dir-2
…
$ svn update --set-depth infinity many-dirs/wanted-dir-3
…
### and so on, and so on, ...
\end{verbatim}

This could be quite tedious, especially since you don\textquotesingle t
even have stubs of these directories in your working copy to deal with.
Such a working copy would also have another characteristic that you
might not expect or desire: if someone else creates any new
subdirectories in this top-level directory, you won\textquotesingle t
receive those when you update your working copy.

Beginning with Subversion 1.6, you can take a different approach. First,
check out the directory in full. Then run
\texttt{svn\ update\ -\/-set-depth\ exclude} on the one subdirectory you
don\textquotesingle t care about.

\begin{verbatim}
$ svn checkout http://svn.example.com/repos/many-dirs
…
$ svn update --set-depth exclude many-dirs/unwanted-dir
D         many-dirs/unwanted-dir
$
\end{verbatim}

This approach leaves your working copy with the same stuff as in the
first approach, but any new subdirectories which appear in the top-level
directory would also show up when you update your working copy. The
downside of this approach is that you have to actually check out that
whole subdirectory that you don\textquotesingle t even want just so you
can tell Subversion that you don\textquotesingle t want it. This might
not even be possible if that subdirectory is too large to fit on your
disk (which might, after all, be the very reason you
don\textquotesingle t want it in your working copy).

While the functionality for excluding an existing item from a working
copy was hung off of the \texttt{svn\ update} command, you might have
noticed that the output from
\texttt{svn\ update\ -\/-set-depth\ exclude} differs from that of a
normal update operation. This output betrays the fact that, under the
hood, exclusion is a completely client-side operation, very much unlike
a typical update.

In such a situation, you might consider a compromise approach. First,
check out the top-level directory with \texttt{-\/-depth\ immediates}.
Then, exclude the directory you don\textquotesingle t want using
\texttt{svn\ update\ -\/-set-depth\ exclude}. Finally, telescope all the
items that remain to infinite depth, which should be fairly easy to do
because they are all addressable in your shell.

\begin{verbatim}
$ svn checkout http://svn.example.com/repos/many-dirs --depth immediates
…
$ svn update --set-depth exclude many-dirs/unwanted-dir
D         many-dirs/unwanted-dir
$ svn update --set-depth infinity many-dirs/*
…
$
\end{verbatim}

Once again, your working copy will have the same stuff as in the
previous two scenarios. But now, any time a new file or subdirectory is
committed to the top-level directory, you\textquotesingle ll receive
it---at an empty depth---when you update your working copy. You can now
decide what to do with such newly appearing working copy items: expand
them into infinite depth, or exclude them altogether.

\section{Locking}\label{svn.advanced.locking}

Subversion\textquotesingle s copy-modify-merge version control model
lives and dies on its data merging algorithms---specifically on how well
those algorithms perform when trying to resolve conflicts caused by
multiple users modifying the same file concurrently. Subversion itself
provides only one such algorithm: a three-way differencing algorithm
that is smart enough to handle data at a granularity of a single line of
text. Subversion also allows you to supplement its content merge
processing with external differencing utilities (as described in
\hyperref[svn.advanced.externaldifftools.diff3]{External diff3} and
\hyperref[svn.advanced.externaldifftools.merge]{External merge}), some
of which may do an even better job, perhaps providing granularity of a
word or a single character of text. But common among those algorithms is
that they generally work only on text files. The landscape starts to
look pretty grim when you start talking about content merges of
nontextual file formats. And when you can\textquotesingle t find a tool
that can handle that type of merging, you begin to run into problems
with the copy-modify-merge model.

Let\textquotesingle s look at a real-life example of where this model
runs aground. Harry and Sally are both graphic designers working on the
same project, a bit of marketing collateral for an automobile mechanic.
Central to the design of a particular poster is an image of a car in
need of some bodywork, stored in a file using the PNG image format. The
poster\textquotesingle s layout is almost finished, and both Harry and
Sally are pleased with the particular photo they chose for their damaged
car---a baby blue 1967 Ford Mustang with an unfortunate bit of crumpling
on the left front fender.

Now, as is common in graphic design work, there\textquotesingle s a
change in plans, which causes the car\textquotesingle s color to be a
concern. So Sally updates her working copy to \texttt{HEAD}, fires up
her photo-editing software, and sets about tweaking the image so that
the car is now cherry red. Meanwhile, Harry, feeling particularly
inspired that day, decides that the image would have greater impact if
the car also appears to have suffered greater impact. He, too, updates
to \texttt{HEAD}, and then draws some cracks on the
vehicle\textquotesingle s windshield. He manages to finish his work
before Sally finishes hers, and after admiring the fruits of his
undeniable talent, he commits the modified image. Shortly thereafter,
Sally is finished with the car\textquotesingle s new finish and tries to
commit her changes. But, as expected, Subversion fails the commit,
informing Sally that her version of the image is now out of date.

Here\textquotesingle s where the difficulty sets in. If Harry and Sally
were making changes to a text file, Sally would simply update her
working copy, receiving Harry\textquotesingle s changes in the process.
In the worst possible case, they would have modified the same region of
the file, and Sally would have to work out by hand the proper resolution
to the conflict. But these aren\textquotesingle t text files---they are
binary images. And while it\textquotesingle s a simple matter to
describe what one would expect the results of this content merge to be,
there is precious little chance that any software exists that is smart
enough to examine the common baseline image that each of these graphic
artists worked against, the changes that Harry made, and the changes
that Sally made, and then spit out an image of a busted-up red Mustang
with a cracked windshield!

Of course, things would have gone more smoothly if Harry and Sally had
serialized their modifications to the image---if, say, Harry had waited
to draw his windshield cracks on Sally\textquotesingle s now-red car, or
if Sally had tweaked the color of a car whose windshield was already
cracked. As is discussed in
\hyperref[svn.basic.vsn-models.copy-merge]{The copy-modify-merge
solution}, most of these types of problems go away entirely where
perfect communication between Harry and Sally exists.\footnote{Communication
  wouldn\textquotesingle t have been such bad medicine for Harry and
  Sally\textquotesingle s Hollywood namesakes, either, for that matter.}
But as one\textquotesingle s version control system is, in fact, one
form of communication, it follows that having that software facilitate
the serialization of nonparallelizable editing efforts is no bad thing.
This is where Subversion\textquotesingle s implementation of the
lock-modify-unlock model steps into the spotlight. This is where we talk
about Subversion\textquotesingle s locking feature, which is similar to
the ``reserved checkouts'' mechanisms of other version control systems.

Subversion\textquotesingle s locking feature exists ultimately to
minimize wasted time and effort. By allowing a user to programmatically
claim the exclusive right to change a file in the repository, that user
can be reasonably confident that any energy he invests on unmergeable
changes won\textquotesingle t be wasted---his commit of those changes
will succeed. Also, because Subversion communicates to other users that
serialization is in effect for a particular versioned object, those
users can reasonably expect that the object is about to be changed by
someone else. They, too, can then avoid wasting their time and energy on
unmergeable changes that won\textquotesingle t be committable due to
eventual out-of-dateness.

When referring to Subversion\textquotesingle s locking feature, one is
actually talking about a fairly diverse collection of behaviors, which
include the ability to lock a versioned file\footnote{Subversion does
  not currently allow locks on directories.} (claiming the exclusive
right to modify the file), to unlock that file (yielding that exclusive
right to modify), to see reports about which files are locked and by
whom, to annotate files for which locking before editing is strongly
advised, and so on. In this section, we\textquotesingle ll cover all of
these facets of the larger locking feature.

\protect\phantomsection\label{svn.advanced.locking.meanings}
The Many Meanings of ``Lock''

In this section, and almost everywhere in this book, the words ``lock''
and ``locking'' describe a mechanism for mutual exclusion between users
to avoid clashing commits. Unfortunately, there are other sorts of
``lock'' with which Subversion, and therefore this book, sometimes needs
to be concerned.

Subversion uses administrative locks internally to prevent clashes
between multiple Subversion clients operating on the same working copy.
This is the sort of lock indicated by an \texttt{L} in the third column
of \texttt{svn\ status} output, and removed by the \texttt{svn\ cleanup}
command, as described in \hyperref[svn.tour.cleanup]{Sometimes You Just
Need to Clean Up}.

Administrators using the older Berkeley DB repository backend will need
to be familiar with database locks, which exist to prevent clashes
between multiple programs trying to access the database. This is the
sort of lock whose unwanted persistence after an error can cause a
repository to be ``wedged,'' as described in
\hyperref[svn.berkeleydb.maintenance.recovery]{Berkeley DB Recovery}.

Additionally, there are svnsync locks, which effect mutual exclusion
between multiple instances of the \texttt{svnsync} command that write to
the same mirror of a repository. This is the sort of lock implemented by
the \texttt{svn:sync-lock} revision property, as described in
\hyperref[svn.reposadmin.maint.replication.svnsync]{Replication with
svnsync}.

Next, there are svnrdump locks. These are very much like svnsync locks,
but are associated with the \texttt{svnrdump\ load} command (described
in \hyperref[svn.reposadmin.maint.migrate.svnrdump]{Repository data
migration using svnrdump}) instead of \texttt{svnsync}.

Finally, there are SQLite locks. These are used by the SQLite library
(\href{https://www.sqlite.org/}{}) to serialize access to SQLite
databases used by Subversion under the hood. See, for example, the
\texttt{exclusive-locking} option in
\hyperref[svn.advanced.confarea.opts.config]{General configuration}.

You can generally forget about these other kinds of locks until
something goes wrong that requires you to care about them. In this book,
``lock'' means the first sort unless the contrary is either clear from
context or explicitly stated.

\subsection{Creating Locks}\label{svn.advanced.locking.creation}

In the Subversion repository, a lock is a piece of metadata that grants
exclusive access to one user to change a file. This user is said to be
the lock owner. Each lock also has a unique identifier, typically a long
string of characters, known as the lock token. The repository manages
locks, ultimately handling their creation, enforcement, and removal. If
any commit transaction attempts to modify or delete a locked file (or
delete one of the parent directories of the file), the repository will
demand two pieces of information---that the client performing the commit
be authenticated as the lock owner, and that the lock token has been
provided as part of the commit process as a form of proof that the
client knows which lock it is using.

To demonstrate lock creation, let\textquotesingle s refer back to our
example of multiple graphic designers working on the same binary image
files. Harry has decided to change a JPEG image. To prevent other people
from committing changes to the file while he is modifying it (as well as
alerting them that he is about to change it), he locks the file in the
repository using the \texttt{svn\ lock} command.

\begin{verbatim}
$ svn lock banana.jpg -m "Editing file for tomorrow's release."
'banana.jpg' locked by user 'harry'.
$
\end{verbatim}

The preceding example demonstrates a number of new things. First, notice
that Harry passed the \texttt{-\/-message} (\texttt{-m}) option to
\texttt{svn\ lock}. Similar to \texttt{svn\ commit}, the
\texttt{svn\ lock} command can take comments---via either
\texttt{-\/-message} (\texttt{-m}) or \texttt{-\/-file}
(\texttt{-F})---to describe the reason for locking the file. Unlike
\texttt{svn\ commit}, however, \texttt{svn\ lock} will not demand a
message by launching your preferred text editor. Lock comments are
optional, but still recommended to aid communication.

Second, the lock attempt succeeded. This means that the file
wasn\textquotesingle t already locked, and that Harry had the latest
version of the file. If Harry\textquotesingle s working copy of the file
had been out of date, the repository would have rejected the request,
forcing Harry to \texttt{svn\ update} and reattempt the locking command.
The locking command would also have failed if the file had already been
locked by someone else.

As you can see, the \texttt{svn\ lock} command prints confirmation of
the successful lock. At this point, the fact that the file is locked
becomes apparent in the output of the \texttt{svn\ status} and
\texttt{svn\ info} reporting subcommands.

\begin{verbatim}
$ svn status
     K  banana.jpg

$ svn info banana.jpg
Path: banana.jpg
Name: banana.jpg
Working Copy Root Path: /home/harry/project
URL: http://svn.example.com/repos/project/banana.jpg
Relative URL: ^/banana.jpg
Repository Root: http://svn.example.com/repos/project
Repository UUID: edb2f264-5ef2-0310-a47a-87b0ce17a8ec
Revision: 2198
Node Kind: file
Schedule: normal
Last Changed Author: frank
Last Changed Rev: 1950
Last Changed Date: 2006-03-15 12:43:04 -0600 (Wed, 15 Mar 2006)
Text Last Updated: 2006-06-08 19:23:07 -0500 (Thu, 08 Jun 2006)
Properties Last Updated: 2006-06-08 19:23:07 -0500 (Thu, 08 Jun 2006)
Checksum: 3b110d3b10638f5d1f4fe0f436a5a2a5
Lock Token: opaquelocktoken:0c0f600b-88f9-0310-9e48-355b44d4a58e
Lock Owner: harry
Lock Created: 2006-06-14 17:20:31 -0500 (Wed, 14 Jun 2006)
Lock Comment (1 line):
Editing file for tomorrow's release.

$
\end{verbatim}

The fact that the \texttt{svn\ info} command, which does not contact the
repository when run against working copy paths, can display the lock
token reveals an important piece of information about those tokens: they
are cached in the working copy. The presence of the lock token is
critical. It gives the working copy authorization to make use of the
lock later on. Also, the \texttt{svn\ status} command shows a \texttt{K}
next to the file (short for locKed), indicating that the lock token is
present.

Regarding Lock Tokens

A lock token isn\textquotesingle t an authentication token, so much as
an \emph{authorization} token. The token isn\textquotesingle t a
protected secret. In fact, a lock\textquotesingle s unique token is
discoverable by anyone who runs \texttt{svn\ info\ URL}. A lock token is
special only when it lives inside a working copy. It\textquotesingle s
proof that the lock was created in that particular working copy, and not
somewhere else by some other client. Merely authenticating as the lock
owner isn\textquotesingle t enough to prevent accidents.

For example, suppose you lock a file using a computer at your office,
but leave work for the day before you finish your changes to that file.
It should not be possible to accidentally commit changes to that same
file from your home computer later that evening simply because
you\textquotesingle ve authenticated as the lock\textquotesingle s
owner. In other words, the lock token prevents one piece of
Subversion-related software from undermining the work of another. (In
our example, if you really need to change the file from an alternative
working copy, you would need to break the lock and relock the file.)

Now that Harry has locked \texttt{banana.jpg}, Sally is unable to change
or delete that file:

\begin{verbatim}
$ svn delete banana.jpg
D         banana.jpg
$ svn commit -m "Delete useless file."
Deleting       banana.jpg
svn: E175002: Commit failed (details follow):
svn: E175002: Server sent unexpected return value (423 Locked) in response to 
DELETE request for '/repos/project/!svn/wrk/64bad3a9-96f9-0310-818a-df4224ddc
35d/banana.jpg'
$
\end{verbatim}

But Harry, after touching up the banana\textquotesingle s shade of
yellow, is able to commit his changes to the file.
That\textquotesingle s because he authenticates as the lock owner and
also because his working copy holds the correct lock token:

\begin{verbatim}
$ svn status
M    K  banana.jpg
$ svn commit -m "Make banana more yellow"
Sending        banana.jpg
Transmitting file data .
Committed revision 2201.
$ svn status
$
\end{verbatim}

Notice that after the commit is finished, \texttt{svn\ status} shows
that the lock token is no longer present in the working copy. This is
the standard behavior of \texttt{svn\ commit}---it searches the working
copy (or list of targets, if you provide such a list) for local
modifications and sends all the lock tokens it encounters during this
walk to the server as part of the commit transaction. After the commit
completes successfully, all of the repository locks that were mentioned
are released---\emph{even on files that weren\textquotesingle t
committed}. This is meant to discourage users from being sloppy about
locking or from holding locks for too long. If Harry haphazardly locks
30 files in a directory named \texttt{images} because
he\textquotesingle s unsure of which files he needs to change, yet
changes only four of those files, when he runs
\texttt{svn\ commit\ images}, the process will still release all 30
locks.

This behavior of automatically releasing locks can be overridden with
the \texttt{-\/-no-unlock} option to \texttt{svn\ commit}. This is best
used for those times when you want to commit changes, but still plan to
make more changes and thus need to retain existing locks. You can also
make this your default behavior by setting the \texttt{no-unlock}
runtime configuration option (see
\hyperref[svn.advanced.confarea]{Runtime Configuration Area}).

Of course, locking a file doesn\textquotesingle t oblige one to commit a
change to it. The lock can be released at any time with a simple
\texttt{svn\ unlock} command:

\begin{verbatim}
$ svn unlock banana.c
'banana.c' unlocked.
\end{verbatim}

\subsection{Discovering Locks}\label{svn.advanced.locking.discovery}

When a commit fails due to someone else\textquotesingle s locks,
it\textquotesingle s fairly easy to learn about them. The easiest way is
to run \texttt{svn\ status\ -u}:

\begin{verbatim}
$ svn status -u
M               23   bar.c
M    O          32   raisin.jpg
        *       72   foo.h
Status against revision:     105
$
\end{verbatim}

In this example, Sally can see not only that her copy of \texttt{foo.h}
is out of date, but also that one of the two modified files she plans to
commit is locked in the repository. The \texttt{O} symbol stands for
``Other,'' meaning that a lock exists on the file and was created by
somebody else. If she were to attempt a commit, the lock on
\texttt{raisin.jpg} would prevent it. Sally is left wondering who made
the lock, when, and why. Once again, \texttt{svn\ info} has the answers:

\begin{verbatim}
$ svn info ^/raisin.jpg
Path: raisin.jpg
Name: raisin.jpg
URL: http://svn.example.com/repos/project/raisin.jpg
Relative URL: ^/raisin.jpg
Repository Root: http://svn.example.com/repos/project
Repository UUID: edb2f264-5ef2-0310-a47a-87b0ce17a8ec
Revision: 105
Node Kind: file
Last Changed Author: sally
Last Changed Rev: 32
Last Changed Date: 2006-01-25 12:43:04 -0600 (Sun, 25 Jan 2006)
Lock Token: opaquelocktoken:fc2b4dee-98f9-0310-abf3-653ff3226e6b
Lock Owner: harry
Lock Created: 2006-02-16 13:29:18 -0500 (Thu, 16 Feb 2006)
Lock Comment (1 line):
Need to make a quick tweak to this image.
$
\end{verbatim}

Just as you can use \texttt{svn\ info} to examine objects in the working
copy, you can also use it to examine objects in the repository. If the
main argument to \texttt{svn\ info} is a working copy path, then all of
the working copy\textquotesingle s cached information is displayed; any
mention of a lock means that the working copy is holding a lock token
(if a file is locked by another user or in another working copy,
\texttt{svn\ info} on a working copy path will show no lock information
at all). If the main argument to \texttt{svn\ info} is a URL, the
information reflects the latest version of an object in the repository,
and any mention of a lock describes the current lock on the object.

So in this particular example, Sally can see that Harry locked the file
on February 16 to ``make a quick tweak.'' It being June, she suspects
that he probably forgot all about the lock. She might phone Harry to
complain and ask him to release the lock. If he\textquotesingle s
unavailable, she might try to forcibly break the lock herself or ask an
administrator to do so.

\subsection{Breaking and Stealing
Locks}\label{svn.advanced.locking.break-steal}

A repository lock isn\textquotesingle t sacred---in
Subversion\textquotesingle s default configuration state, locks can be
released not only by the person who created them, but by anyone. When
somebody other than the original lock creator destroys a lock, we refer
to this as breaking the lock.

From the administrator\textquotesingle s chair, it\textquotesingle s
simple to break locks. The \texttt{svnlook} and \texttt{svnadmin}
programs have the ability to display and remove locks directly from the
repository. (For more information about these tools, see
\hyperref[svn.reposadmin.maint.tk]{An Administrator\textquotesingle s
Toolkit}.)

\begin{verbatim}
$ svnadmin lslocks /var/svn/repos
Path: /project2/images/banana.jpg
UUID Token: opaquelocktoken:c32b4d88-e8fb-2310-abb3-153ff1236923
Owner: frank
Created: 2006-06-15 13:29:18 -0500 (Thu, 15 Jun 2006)
Expires: 
Comment (1 line):
Still improving the yellow color.

Path: /project/raisin.jpg
UUID Token: opaquelocktoken:fc2b4dee-98f9-0310-abf3-653ff3226e6b
Owner: harry
Created: 2006-02-16 13:29:18 -0500 (Thu, 16 Feb 2006)
Expires: 
Comment (1 line):
Need to make a quick tweak to this image.

$ svnadmin rmlocks /var/svn/repos /project/raisin.jpg
Removed lock on '/project/raisin.jpg'.
$
\end{verbatim}

The more interesting option is to allow users to break each
other\textquotesingle s locks over the network. To do this, Sally simply
needs to pass the \texttt{-\/-force} option to the \texttt{svn\ unlock}
command:

\begin{verbatim}
$ svn status -u
M               23   bar.c
M    O          32   raisin.jpg
        *       72   foo.h
Status against revision:     105
$ svn unlock raisin.jpg
svn: E195013: 'raisin.jpg' is not locked in this working copy
$ svn info raisin.jpg | grep ^URL
URL: http://svn.example.com/repos/project/raisin.jpg
$ svn unlock http://svn.example.com/repos/project/raisin.jpg
svn: warning: W160039: Unlock failed on 'raisin.jpg' (403 Forbidden)
$ svn unlock --force http://svn.example.com/repos/project/raisin.jpg
'raisin.jpg' unlocked.
$
\end{verbatim}

Now, Sally\textquotesingle s initial attempt to unlock failed because
she ran \texttt{svn\ unlock} directly on her working copy of the file,
and no lock token was present. To remove the lock directly from the
repository, she needs to pass a URL to \texttt{svn\ unlock}. Her first
attempt to unlock the URL fails, because she can\textquotesingle t
authenticate as the lock owner (nor does she have the lock token). But
when she passes \texttt{-\/-force}, the authentication and authorization
requirements are ignored, and the remote lock is broken.

Simply breaking a lock may not be enough. In the running example, Sally
may not only want to break Harry\textquotesingle s long-forgotten lock,
but relock the file for her own use. She can accomplish this by using
\texttt{svn\ unlock} with \texttt{-\/-force} and then \texttt{svn\ lock}
back-to-back, but there\textquotesingle s a small chance that somebody
else might lock the file between the two commands. The simpler thing to
do is to steal the lock, which involves breaking and relocking the file
all in one atomic step. To do this, Sally passes the \texttt{-\/-force}
option to \texttt{svn\ lock}:

\begin{verbatim}
$ svn lock raisin.jpg
svn: warning: W160035: Path '/project/raisin.jpg' is already locked by user 'h
arry' in filesystem '/var/svn/repos/db'
$ svn lock --force raisin.jpg
'raisin.jpg' locked by user 'sally'.
$
\end{verbatim}

In any case, whether the lock is broken or stolen, Harry may be in for a
surprise. Harry\textquotesingle s working copy still contains the
original lock token, but that lock no longer exists. The lock token is
said to be defunct. The lock represented by the lock token has either
been broken (no longer in the repository) or stolen (replaced with a
different lock). Either way, Harry can see this by asking
\texttt{svn\ status} to contact the repository:

\begin{verbatim}
$ svn status
     K  raisin.jpg
$ svn status -u
     B          32   raisin.jpg
Status against revision:     105
$ svn update
Updating '.':
  B  raisin.jpg
Updated to revision 105.
$ svn status
$
\end{verbatim}

If the repository lock was broken, then
\texttt{svn\ status\ -\/-show-updates} (\texttt{-u}) displays a
\texttt{B} (Broken) symbol next to the file. If a new lock exists in
place of the old one, then a \texttt{T} (sTolen) symbol is shown.
Finally, \texttt{svn\ update} notices any defunct lock tokens and
removes them from the working copy.

Locking Policies

Different systems have different notions of how strict a lock should be.
Some folks argue that locks must be strictly enforced at all costs,
releasable only by the original creator or administrator. They argue
that if anyone can break a lock, chaos runs rampant and the whole point
of locking is defeated. The other side argues that locks are first and
foremost a communication tool. If users are constantly breaking each
other\textquotesingle s locks, it represents a cultural failure within
the team and the problem falls outside the scope of software
enforcement.

Subversion defaults to the ``softer'' approach, but still allows
administrators to create stricter enforcement policies through the use
of hook scripts. In particular, the pre-lock and pre-unlock hooks allow
administrators to decide when lock creation and lock releases are
allowed to happen. Depending on whether a lock already exists, these two
hooks can decide whether to allow a certain user to break or steal a
lock. The post-lock and post-unlock hooks are also available, and can be
used to send email after locking actions. To learn more about repository
hooks, see \hyperref[svn.reposadmin.hooks]{Implementing Repository
Hooks}.

\subsection{Lock
Communication}\label{svn.advanced.locking.lock-communication}

We\textquotesingle ve seen how \texttt{svn\ lock} and
\texttt{svn\ unlock} can be used to create, release, break, and steal
locks. This satisfies the goal of serializing commit access to a file.
But what about the larger problem of preventing wasted time?

For example, suppose Harry locks an image file and then begins editing
it. Meanwhile, miles away, Sally wants to do the same thing. She
doesn\textquotesingle t think to run \texttt{svn\ status\ -u}, so she
has no idea that Harry has already locked the file. She spends hours
editing the file, and when she tries to commit her change, she discovers
that either the file is locked or that she\textquotesingle s out of
date. Regardless, her changes aren\textquotesingle t mergeable with
Harry\textquotesingle s. One of these two people has to throw away his
or her work, and a lot of time has been wasted.

Subversion\textquotesingle s solution to this problem is to provide a
mechanism to remind users that a file ought to be locked \emph{before}
the editing begins. The mechanism is a special property:
\texttt{svn:needs-lock}. If that property is attached to a file
(regardless of its value, which is irrelevant), Subversion will try to
use filesystem-level permissions to make the file read-only---unless, of
course, the user has explicitly locked the file. When a lock token is
present (as a result of using \texttt{svn\ lock}), the file becomes
read/write. When the lock is released, the file becomes read-only again.

The theory, then, is that if the image file has this property attached,
Sally would immediately notice something is strange when she opens the
file for editing: many applications alert users immediately when a
read-only file is opened for editing, and nearly all would prevent her
from saving changes to the file. This reminds her to lock the file
before editing, whereby she discovers the preexisting lock:

\begin{verbatim}
$ /usr/local/bin/gimp raisin.jpg
gimp: error: file is read-only!
$ ls -l raisin.jpg
-r--r--r--   1 sally   sally   215589 Jun  8 19:23 raisin.jpg
$ svn lock raisin.jpg
svn: warning: W160035: Path '/project/raisin.jpg' is already locked by user 'h
arry' in filesystem '/var/svn/repos/db'
$ svn info http://svn.example.com/repos/project/raisin.jpg | grep Lock
Lock Token: opaquelocktoken:fc2b4dee-98f9-0310-abf3-653ff3226e6b
Lock Owner: harry
Lock Created: 2006-06-08 07:29:18 -0500 (Thu, 08 June 2006)
Lock Comment (1 line):
Making some tweaks.  Locking for the next two hours.
$
\end{verbatim}

Users and administrators alike are encouraged to attach the
\texttt{svn:needs-lock} property to any file that cannot be contextually
merged. This is the primary technique for encouraging good locking
habits and preventing wasted effort.

Note that this property is a communication tool that works independently
from the locking system. In other words, any file can be locked, whether
or not this property is present. And conversely, the presence of this
property doesn\textquotesingle t make the repository require a lock when
committing.

Unfortunately, the system isn\textquotesingle t flawless.
It\textquotesingle s possible that even when a file has the property,
the read-only reminder won\textquotesingle t always work. Sometimes
applications misbehave and ``hijack'' the read-only file, silently
allowing users to edit and save the file anyway. There\textquotesingle s
not much that Subversion can do in this situation---at the end of the
day, there\textquotesingle s simply no substitution for good
interpersonal communication.\footnote{Except, perhaps, a classic Vulcan
  mind-meld.}

\section{Externals Definitions}\label{svn.advanced.externals}

Sometimes it is useful to construct a working copy that is made out of a
number of different checkouts. For example, you may want different
subdirectories to come from different locations in a repository or
perhaps from different repositories altogether. You could certainly set
up such a scenario by hand---using \texttt{svn\ checkout} to create the
sort of nested working copy structure you are trying to achieve. But if
this layout is important for everyone who uses your repository, every
other user will need to perform the same checkout operations that you
did.

Fortunately, Subversion provides support for externals definitions. An
externals definition is a mapping of a local directory to the URL---and
ideally a particular revision---of a versioned directory. In Subversion,
you declare externals definitions in groups using the
\texttt{svn:externals} property. You can create or modify this property
using \texttt{svn\ propset} or \texttt{svn\ propedit} (see
\hyperref[svn.advanced.props.manip]{Manipulating Properties}). It can be
set on any versioned directory, and its value describes both the
external repository location and the client-side directory to which that
location should be checked out.

The convenience of the \texttt{svn:externals} property is that once it
is set on a versioned directory, everyone who checks out a working copy
with that directory also gets the benefit of the externals definition.
In other words, once one person has made the effort to define the nested
working copy structure, no one else has to bother---Subversion will,
after checking out the original working copy, automatically also check
out the external working copies.

The relative target subdirectories of externals definitions \emph{must
not} already exist on your or other users\textquotesingle{}
systems---Subversion will create them when it checks out the external
working copy.

You also get in the externals definition design all the regular benefits
of Subversion properties. The definitions are versioned. If you need to
change an externals definition, you can do so using the regular property
modification subcommands. When you commit a change to the
\texttt{svn:externals} property, Subversion will synchronize the
checked-out items against the changed externals definition when you next
run \texttt{svn\ update}. The same thing will happen when others update
their working copies and receive your changes to the externals
definition.

Because the \texttt{svn:externals} property has a multiline value, we
strongly recommend that you use \texttt{svn\ propedit} instead of
\texttt{svn\ propset}.

Subversion releases prior to 1.5 honor an externals definition format
that is a multiline table of subdirectories (relative to the versioned
directory on which the property is set), optional revision flags, and
fully qualified, absolute Subversion repository URLs. An example of this
might look as follows:

\begin{verbatim}
$ svn propget svn:externals calc
# Resources versioned elsewhere.
third-party/sounds             http://svn.example.com/repos/sounds
third-party/skins -r148        http://svn.example.com/skinproj
third-party/skins/toolkit -r21 http://svn.example.com/skin-maker
$
\end{verbatim}

In the previous example, three externals have been defined. The first is
``unpinned''---it effectively refers to the HEAD revision of its target.
The other two have specific revisions declared. Also demonstrated is the
fact that lines that begin with a hash (\texttt{\#}) are treated as
comments and ignored by the externals definitions parsing logic.

When someone checks out a working copy of the \texttt{calc} directory
referred to in the previous example, Subversion also continues to check
out the items found in its externals definition.

\begin{verbatim}
$ svn checkout http://svn.example.com/repos/calc
A    calc
A    calc/Makefile
A    calc/integer.c
A    calc/button.c
Checked out revision 148.

Fetching external item into calc/third-party/sounds
A    calc/third-party/sounds/ding.ogg
A    calc/third-party/sounds/dong.ogg
A    calc/third-party/sounds/clang.ogg
…
A    calc/third-party/sounds/bang.ogg
A    calc/third-party/sounds/twang.ogg
Checked out revision 14.

Fetching external item into calc/third-party/skins
…
\end{verbatim}

As of Subversion 1.5, though, a new format of the \texttt{svn:externals}
property is supported. Externals definitions are still multiline, but
the order and format of the various pieces of information have changed.
The new syntax more closely mimics the order of arguments you might pass
to \texttt{svn\ checkout}: the optional revision flags come first, then
the external Subversion repository URL, and finally the relative local
subdirectory. Notice, though, that this time we didn\textquotesingle t
say ``fully qualified, absolute Subversion repository URLs.''
That\textquotesingle s because the new format supports relative URLs and
URLs that carry peg revisions. The previous example of an externals
definition might, in Subversion 1.5, look like the following:

\begin{verbatim}
$ svn propget svn:externals calc
# Resources versioned elsewhere.
      http://svn.example.com/repos/sounds third-party/sounds
-r148 http://svn.example.com/skinproj third-party/skins
-r21  http://svn.example.com/skin-maker third-party/skins/toolkit
$
\end{verbatim}

Or, making use of the peg revision syntax (which we describe in detail
in \hyperref[svn.advanced.pegrevs]{Peg and Operative Revisions}), it
might appear as:

\begin{verbatim}
$ svn propget svn:externals calc
# Resources versioned elsewhere.
http://svn.example.com/repos/sounds third-party/sounds
http://svn.example.com/skinproj@148 third-party/skins
http://svn.example.com/skin-maker@21 third-party/skins/toolkit
$
\end{verbatim}

You should seriously consider using explicit revision numbers in all of
your externals definitions. Doing so means that you get to decide when
to pull down a different snapshot of external information, and exactly
which snapshot to pull. Besides avoiding the surprise of getting changes
to third-party repositories that you might not have any control over,
using explicit revision numbers also means that as you backdate your
working copy to a previous revision, your externals definitions will
also revert to the way they looked in that previous revision, which in
turn means that the external working copies will be updated to match the
way \emph{they} looked back when your repository was at that previous
revision. For software projects, this could be the difference between a
successful and a failed build of an older snapshot of your complex
codebase.

For most repositories, these three ways of formatting the externals
definitions have the same ultimate effect. They all bring the same
benefits. Unfortunately, they all bring the same annoyances, too. Since
the definitions shown use absolute URLs, moving or copying a directory
to which they are attached will not affect what gets checked out as an
external (though the relative local target subdirectory will, of course,
move with the renamed directory). This can be confusing---even
frustrating---in certain situations. For example, say you have a
top-level directory named \texttt{my-project}, and
you\textquotesingle ve created an externals definition on one of its
subdirectories (\texttt{my-project/some-dir}) that tracks the latest
revision of another of its subdirectories
(\texttt{my-project/external-dir}).

\begin{verbatim}
$ svn checkout http://svn.example.com/projects .
A    my-project
A    my-project/some-dir
A    my-project/external-dir
…
Fetching external item into 'my-project/some-dir/subdir'
Checked out external at revision 11.

Checked out revision 11.
$ svn propget svn:externals my-project/some-dir
subdir http://svn.example.com/projects/my-project/external-dir

$
\end{verbatim}

Now you use \texttt{svn\ move} to rename the \texttt{my-project}
directory. At this point, your externals definition will still refer to
a path under the \texttt{my-project} directory, even though that
directory no longer exists.

\begin{verbatim}
$ svn move -q my-project renamed-project
$ svn commit -m "Rename my-project to renamed-project."
Deleting       my-project
Adding         renamed-project

Committed revision 12.
$ svn update
Updating '.':

svn: warning: W200000: Error handling externals definition for 'renamed-projec
t/some-dir/subdir':
svn: warning: W170000: URL 'http://svn.example.com/projects/my-project/externa
l-dir' at revision 12 doesn't exist
At revision 12.
svn: E205011: Failure occurred processing one or more externals definitions
$
\end{verbatim}

Also, absolute URLs can cause problems with repositories that are
available via multiple URL schemes. For example, if your Subversion
server is configured to allow everyone to check out the repository over
\texttt{http://} or \texttt{https://}, but only allow commits to come in
via \texttt{https://}, you have an interesting problem on your hands. If
your externals definitions use the \texttt{http://} form of the
repository URLs, you won\textquotesingle t be able to commit anything
from the working copies created by those externals. On the other hand,
if they use the \texttt{https://} form of the URLs, anyone who might be
checking out via \texttt{http://} because his client
doesn\textquotesingle t support \texttt{https://} will be unable to
fetch the external items. Be aware, too, that if you need to reparent
your working copy (using \texttt{svn\ relocate}), externals definitions
will \emph{not} also be reparented.

Subversion 1.5 takes a huge step in relieving these frustrations. As
mentioned earlier, the URLs used in the new externals definition format
can be relative, and Subversion provides syntax magic for specifying
multiple flavors of URL relativity.

\protect\phantomsection\label{svn.advanced.externals.urlmagic}{}

\begin{description}
\item[\texttt{../}]
Relative to the URL of the directory on which the \texttt{svn:externals}
property is set
\item[\texttt{\^{}/}]
Relative to the root of the repository in which the
\texttt{svn:externals} property is versioned
\item[\texttt{//}]
Relative to the scheme of the URL of the directory on which the
\texttt{svn:externals} property is set
\item[\texttt{/}]
Relative to the root URL of the server on which the
\texttt{svn:externals} property is versioned
\item[\texttt{\^{}/../REPO-NAME}]
Relative to a sibling repository beneath the same \texttt{SVNParentPath}
location as the repository in which the \texttt{svn:externals} is
defined.
\end{description}

So, looking a fourth time at our previous externals definition example,
and making use of the new absolute URL syntax in various ways, we might
now see:

\begin{verbatim}
$ svn propget svn:externals calc
# Resources versioned elsewhere.
^/sounds third-party/sounds
/skinproj@148 third-party/skins
//svn.example.com/skin-maker@21 third-party/skins/toolkit
$
\end{verbatim}

Subversion 1.6 brought two more improvements to externals definitions.
First, it added a quoting and escape mechanism to the syntax so that the
path of the external working copy may contain whitespace. This was
previously problematic, of course, because whitespace is used to delimit
the fields in an externals definition. Now you need only wrap such a
path specification in double-quote (\texttt{"}) characters or escape the
problematic characters in the path with a backslash
(\texttt{\textbackslash{}}) character. Of course, if you have spaces in
the \emph{URL} portion of the external definition, you should use the
standard URI-encoding mechanism to represent those.

\begin{verbatim}
$ svn propget svn:externals paint
http://svn.thirdparty.com/repos/My%20Project "My Project"
http://svn.thirdparty.com/repos/%22Quotes%20Too%22 \"Quotes\ Too\"
$
\end{verbatim}

Subversion 1.6 also introduced support for external definitions for
files. File externals are configured just like externals for directories
and appear as a versioned file in the working copy.

For example, let\textquotesingle s say you had the file
\texttt{/trunk/bikeshed/blue.html} in your repository, and you wanted
this file, as it appeared in revision 40, to appear in your working copy
of \texttt{/trunk/www/} as \texttt{green.html}.

The externals definition required to achieve this should look familiar
by now:

\begin{verbatim}
$ svn propget svn:externals www/
^/trunk/bikeshed/blue.html@40 green.html
$ svn update
Updating '.':

Fetching external item into 'www'
E    www/green.html
Updated external to revision 40.

Update to revision 103.
$ svn status
    X   www/green.html
$
\end{verbatim}

As you can see in the previous output, Subversion denotes file externals
with the letter \texttt{E} when they are fetched into the working copy,
and with the letter \texttt{X} when showing the working copy status.

While directory externals can place the external directory at any depth,
and any missing intermediate directories will be created, file externals
must be placed into a working copy that is already checked out.

When examining the file external with \texttt{svn\ info}, you can see
the URL and revision the external is coming from:

\begin{verbatim}
$ svn info www/green.html 
Path: www/green.html
Name: green.html
Working Copy Root Path: /home/harry/projects/my-project
URL: http://svn.example.com/projects/my-project/trunk/bikeshed/blue.html
Relative URL: ^/trunk/bikeshed/blue.html
Repository Root: http://svn.example.com/projects/my-project
Repository UUID: b2a368dc-7564-11de-bb2b-113435390e17
Revision: 40
Node kind: file
Schedule: normal
Last Changed Author: harry
Last Changed Rev: 40
Last Changed Date: 2009-07-20 20:38:20 +0100 (Mon, 20 Jul 2009)
Text Last Updated: 2009-07-20 23:22:36 +0100 (Mon, 20 Jul 2009)
Checksum: 01a58b04617b92492d99662c3837b33b
$
\end{verbatim}

Because file externals appear in the working copy as versioned files,
they can be modified and even committed if they reference a file at the
HEAD revision. The committed changes will then appear in the external as
well as the file referenced by the external. However, in our example, we
pinned the external to an older revision, so attempting to commit the
external fails:

\begin{verbatim}
$ svn status
M   X   www/green.html
$ svn commit -m "change the color" www/green.html
Sending        www/green.html
svn: E155011: Commit failed (details follow):
svn: E155011: File '/trunk/bikeshed/blue.html' is out of date
$
\end{verbatim}

Keep this in mind when defining file externals. If you need the external
to refer to a certain revision of a file you will not be able to modify
the external. If you want to be able to modify the external, you cannot
specify a revision other than the \texttt{HEAD} revision, which is
implied if no revision is specified.

Unfortunately, the support which exists for externals definitions in
Subversion remains less than ideal. Both file and directory externals
have shortcomings. For either type of external, the local subdirectory
part of the definition cannot contain \texttt{..} parent directory
indicators (such as \texttt{../../skins/myskin}). File externals cannot
refer to files from other repositories. A file
external\textquotesingle s URL must always be in the same repository as
the URL that the file external will be inserted into. Also, file
externals cannot be moved or deleted. The \texttt{svn:externals}
property must be modified instead. However, file externals can be
copied.

Perhaps most disappointingly, the working copies created via the
externals definition support are still disconnected from the primary
working copy (on whose versioned directories the \texttt{svn:externals}
property was actually set). And Subversion still truly operates only on
nondisjoint working copies. So, for example, if you want to commit
changes that you\textquotesingle ve made in one or more of those
external working copies, you must run \texttt{svn\ commit} explicitly on
those working copies---committing on the primary working copy will not
recurse into any external ones.

We\textquotesingle ve already mentioned some of the additional
shortcomings of the old \texttt{svn:externals} format and how the newer
Subversion 1.5 format improves upon it. But be careful when making use
of the new format that you don\textquotesingle t inadvertently introduce
new problems. For example, while the latest clients will continue to
recognize and support the original externals definition format, pre-1.5
clients will \emph{not} be able to correctly parse the new format. If
you change all your externals definitions to the newer format, you
effectively force everyone who uses those externals to upgrade their
Subversion clients to a version that can parse them. Also, be careful to
avoid naively relocating the \texttt{-rNNN} portion of the
definition---the older format uses that revision as a peg revision, but
the newer format uses it as an operative revision (with a peg revision
of \texttt{HEAD} unless otherwise specified; see
\hyperref[svn.advanced.pegrevs]{Peg and Operative Revisions} for a full
explanation of the distinction here).

External working copies are still completely self-sufficient working
copies. You can operate directly on them as you would any other working
copy. This can be a handy feature, allowing you to examine an external
working copy independently of any primary working copy whose
\texttt{svn:externals} property caused its instantiation. Be careful,
though, that you don\textquotesingle t inadvertently modify your
external working copy in subtle ways that cause problems. For example,
while an externals definition might specify that the external working
copy should be held at a particular revision number, if you run
\texttt{svn\ update} directly on the external working copy, Subversion
will oblige, and now your external working copy is out of sync with its
declaration in the primary working copy. Using \texttt{svn\ switch} to
directly switch the external working copy (or some portion thereof) to
another URL could cause similar problems if the contents of the primary
working copy are expecting particular contents in the external content.

Besides the \texttt{svn\ checkout}, \texttt{svn\ update},
\texttt{svn\ switch}, and \texttt{svn\ export} commands which actually
manage the disjoint (or disconnected) subdirectories into which
externals are checked out, the \texttt{svn\ status} command also
recognizes externals definitions. It displays a status code of
\texttt{X} for the disjoint external subdirectories, and then recurses
into those subdirectories to display the status of the external items
themselves. You can pass the \texttt{-\/-ignore-externals} option to any
of these subcommands to disable externals definition processing.

\section{Changelists}\label{svn.advanced.changelists}

It is commonplace for a developer to find himself working at any given
time on multiple different, distinct changes to a particular bit of
source code. This isn\textquotesingle t necessarily due to poor planning
or some form of digital masochism. A software engineer often spots bugs
in his peripheral vision while working on some nearby chunk of source
code. Or perhaps he\textquotesingle s halfway through some large change
when he realizes the solution he\textquotesingle s working on is best
committed as several smaller logical units. Often, these logical units
aren\textquotesingle t nicely contained in some module, safely separated
from other changes. The units might overlap, modifying different files
in the same module, or even modifying different lines in the same file.

Developers can employ various work methodologies to keep these logical
changes organized. Some use separate working copies of the same
repository to hold each individual change in progress. Others might
choose to create short-lived feature branches in the repository and use
a single working copy that is constantly switched to point to one such
branch or another. Still others use \texttt{diff} and \texttt{patch}
tools to back up and restore uncommitted changes to and from patch files
associated with each change. Each of these methods has its pros and
cons, and to a large degree, the details of the changes being made
heavily influence the methodology used to distinguish them.

Subversion provides a changelists feature that adds yet another method
to the mix. Changelists are basically arbitrary labels (currently at
most one per file) applied to working copy files for the express purpose
of associating multiple files together. Users of many of
Google\textquotesingle s software offerings are familiar with this
concept already. For example, \href{https://mail.google.com/}{Gmail}
doesn\textquotesingle t provide the traditional folders-based email
organization mechanism. In Gmail, you apply arbitrary labels to emails,
and multiple emails can be said to be part of the same group if they
happen to share a particular label. Viewing only a group of similarly
labeled emails then becomes a simple user interface trick. Many other
Web 2.0 sites have similar mechanisms---consider the ``tags'' used by
sites such as \href{https://www.youtube.com/}{YouTube} and
\href{https://www.flickr.com/}{Flickr}, ``categories'' applied to blog
posts, and so on. Folks understand today that organization of data is
critical, but that how that data is organized needs to be a flexible
concept. The old files-and-folders paradigm is too rigid for some
applications.

Subversion\textquotesingle s changelist support allows you to create
changelists by applying labels to files you want to be associated with
that changelist, remove those labels, and limit the scope of the files
on which its subcommands operate to only those bearing a particular
label. In this section, we\textquotesingle ll look in detail at how to
do these things.

\subsection{Creating and Modifying
Changelists}\label{svn.advanced.changelists.creating}

You can create, modify, and delete changelists using the
\texttt{svn\ changelist} command. More accurately, you use this command
to set or unset the changelist association of a particular working copy
file. A changelist is effectively created the first time you label a
file with that changelist; it is deleted when you remove that label from
the last file that had it. Let\textquotesingle s examine a usage
scenario that demonstrates these concepts.

Harry is fixing some bugs in the calculator
application\textquotesingle s mathematics logic. His work leads him to
change a couple of files:

\begin{verbatim}
$ svn status
M       integer.c
M       mathops.c
$
\end{verbatim}

While testing his bug fix, Harry notices that his changes bring to light
a tangentially related bug in the user interface logic found in
\texttt{button.c}. Harry decides that he\textquotesingle ll go ahead and
fix that bug, too, as a separate commit from his math fixes. Now, in a
small working copy with only a handful of files and few logical changes,
Harry can probably keep his two logical change groupings mentally
organized without any problem. But today he\textquotesingle s going to
use Subversion\textquotesingle s changelists feature as a special favor
to the authors of this book.

Harry first creates a changelist and associates with it the two files
he\textquotesingle s already changed. He does this by using the
\texttt{svn\ changelist} command to assign the same arbitrary changelist
name to those files:

\begin{verbatim}
$ svn changelist math-fixes integer.c mathops.c
A [math-fixes] integer.c
A [math-fixes] mathops.c
$ svn status

--- Changelist 'math-fixes':
M       integer.c
M       mathops.c
$
\end{verbatim}

As you can see, the output of \texttt{svn\ status} reflects this new
grouping.

Harry now sets off to fix the secondary UI problem. Since he knows which
file he\textquotesingle ll be changing, he assigns that path to a
changelist, too. Unfortunately, Harry carelessly assigns this third file
to the same changelist as the previous two files:

\begin{verbatim}
$ svn changelist math-fixes button.c
A [math-fixes] button.c
$ svn status

--- Changelist 'math-fixes':
        button.c
M       integer.c
M       mathops.c
$
\end{verbatim}

Fortunately, Harry catches his mistake. At this point, he has two
options. He can remove the changelist association from
\texttt{button.c}, and then assign a different changelist name:

\begin{verbatim}
$ svn changelist --remove button.c
D [math-fixes] button.c
$ svn changelist ui-fix button.c
A [ui-fix] button.c
$
\end{verbatim}

Or, he can skip the removal and just assign a new changelist name. In
this case, Subversion will first warn Harry that \texttt{button.c} is
being removed from the first changelist:

\begin{verbatim}
$ svn changelist ui-fix button.c
D [math-fixes] button.c
A [ui-fix] button.c
$ svn status

--- Changelist 'ui-fix':
        button.c

--- Changelist 'math-fixes':
M       integer.c
M       mathops.c
$
\end{verbatim}

Harry now has two distinct changelists present in his working copy, and
\texttt{svn\ status} will group its output according to these changelist
determinations. Notice that even though Harry hasn\textquotesingle t yet
modified \texttt{button.c}, it still shows up in the output of
\texttt{svn\ status} as interesting because it has a changelist
assignment. Changelists can be added to and removed from files at any
time, regardless of whether they contain local modifications.

Harry now fixes the user interface problem in \texttt{button.c}.

\begin{verbatim}
$ svn status

--- Changelist 'ui-fix':
M       button.c

--- Changelist 'math-fixes':
M       integer.c
M       mathops.c
$
\end{verbatim}

\subsection{Changelists As Operation
Filters}\label{svn.advanced.changelists.asfilters}

The visual grouping that Harry sees in the output of
\texttt{svn\ status} as shown in our previous section is nice, but not
entirely useful. The \texttt{status} command is but one of many
operations that he might wish to perform on his working copy.
Fortunately, many of Subversion\textquotesingle s other operations
understand how to operate on changelists via the use of the
\texttt{-\/-changelist} option.

When provided with a \texttt{-\/-changelist} option, Subversion commands
will limit the scope of their operation to only those files to which a
particular changelist name is assigned. If Harry now wants to see the
actual changes he\textquotesingle s made to the files in his
\texttt{math-fixes} changelist, he \emph{could} explicitly list only the
files that make up that changelist on the \texttt{svn\ diff} command
line.

\begin{verbatim}
$ svn diff integer.c mathops.c
Index: integer.c
===================================================================
--- integer.c   (revision 1157)
+++ integer.c   (working copy)
…
Index: mathops.c
===================================================================
--- mathops.c   (revision 1157)
+++ mathops.c   (working copy)
…
$
\end{verbatim}

That works okay for a few files, but what if Harry\textquotesingle s
change touched 20 or 30 files? That would be an annoyingly long list of
explicitly named files. Now that he\textquotesingle s using changelists,
though, Harry can avoid explicitly listing the set of files in his
changelist from now on, and instead provide just the changelist name:

\begin{verbatim}
$ svn diff --changelist math-fixes
Index: integer.c
===================================================================
--- integer.c   (revision 1157)
+++ integer.c   (working copy)
…
Index: mathops.c
===================================================================
--- mathops.c   (revision 1157)
+++ mathops.c   (working copy)
…
$
\end{verbatim}

And when it\textquotesingle s time to commit, Harry can again use the
\texttt{-\/-changelist} option to limit the scope of the commit to files
in a certain changelist. He might commit his user interface fix by doing
the following:

\begin{verbatim}
$ svn commit -m "Fix a UI bug found while working on math logic." \
             --changelist ui-fix
Sending        button.c
Transmitting file data .
Committed revision 1158.
$
\end{verbatim}

In fact, the \texttt{svn\ commit} command provides a second
changelists-related option: \texttt{-\/-keep-changelists}. Normally,
changelist assignments are removed from files after they are committed.
But if \texttt{-\/-keep-changelists} is provided, Subversion will leave
the changelist assignment on the committed (and now unmodified) files.
In any case, committing files assigned to one changelist leaves other
changelists undisturbed.

\begin{verbatim}
$ svn status

--- Changelist 'math-fixes':
M       integer.c
M       mathops.c
$
\end{verbatim}

The \texttt{-\/-changelist} option acts only as a filter for Subversion
command targets, and will not add targets to an operation. For example,
on a commit operation specified as \texttt{svn\ commit\ /path/to/dir},
the target is the directory \texttt{/path/to/dir} and its children (to
infinite depth). If you then add a changelist specifier to that command,
only those files in and under \texttt{/path/to/dir} that are assigned
that changelist name will be considered as targets of the commit---the
commit will not include files located elsewhere (such as in
\texttt{/path/to/another-dir}), regardless of their changelist
assignment, even if they are part of the same working copy as the
operation\textquotesingle s target(s).

Even the \texttt{svn\ changelist} command accepts the
\texttt{-\/-changelist} option. This allows you to quickly and easily
rename or remove a changelist:

\begin{verbatim}
$ svn changelist math-bugs --changelist math-fixes --depth infinity .
D [math-fixes] integer.c
A [math-bugs] integer.c
D [math-fixes] mathops.c
A [math-bugs] mathops.c
$ svn changelist --remove --changelist math-bugs --depth infinity .
D [math-bugs] integer.c
D [math-bugs] mathops.c
$
\end{verbatim}

Finally, you can specify multiple instances of the
\texttt{-\/-changelist} option on a single command line. Doing so limits
the operation you are performing to files found in any of the specified
changesets.

\subsection{Changelist
Limitations}\label{svn.advanced.changelists.limitations}

Subversion\textquotesingle s changelist feature is a handy tool for
grouping working copy files, but it does have a few limitations.
Changelists are artifacts of a particular working copy, which means that
changelist assignments cannot be propagated to the repository or
otherwise shared with other users. Changelists can be assigned only to
files---Subversion doesn\textquotesingle t currently support the use of
changelists with directories. Finally, you can have at most one
changelist assignment on a given working copy file. Here is where the
blog post category and photo service tag analogies break down---if you
find yourself needing to assign a file to multiple changelists,
you\textquotesingle re out of luck.

\section{Network Model}\label{svn.serverconfig.netmodel}

At some point, you\textquotesingle re going to need to understand how
your Subversion client communicates with its server.
Subversion\textquotesingle s networking layer is abstracted, meaning
that Subversion clients exhibit the same general behaviors no matter
what sort of server they are operating against. Whether speaking the
HTTP protocol (\texttt{http://}) with the Apache HTTP Server or speaking
the custom Subversion protocol (\texttt{svn://}) with \texttt{svnserve},
the basic network model is the same. In this section,
we\textquotesingle ll explain the basics of that network model,
including how Subversion manages authentication and authorization
matters.

\subsection{Requests and
Responses}\label{svn.serverconfig.netmodel.reqresp}

The Subversion client spends most of its time managing working copies.
When it needs information from a remote repository, however, it makes a
network request, and the server responds with an appropriate answer. The
details of the network protocol are hidden from the user---the client
attempts to access a URL, and depending on the URL scheme, a particular
protocol is used to contact the server (see
\hyperref[svn.advanced.reposurls]{Addressing the Repository}).

Run \texttt{svn\ -\/-version} to see which URL schemes and protocols the
client knows how to use.

When the server process receives a client request, it often demands that
the client identify itself. It issues an authentication challenge to the
client, and the client responds by providing credentials back to the
server. Once authentication is complete, the server responds with the
original information that the client asked for. Notice that this system
is different from systems such as CVS, where the client preemptively
offers credentials (``logs in'') to the server before ever making a
request. In Subversion, the server ``pulls'' credentials by challenging
the client at the appropriate moment, rather than the client ``pushing''
them. This makes certain operations more elegant. For example, if a
server is configured to allow anyone in the world to read a repository,
the server will never issue an authentication challenge when a client
attempts to \texttt{svn\ checkout}.

If the particular network requests issued by the client result in a new
revision being created in the repository (e.g., \texttt{svn\ commit}),
Subversion uses the authenticated username associated with those
requests as the author of the revision. That is, the authenticated
user\textquotesingle s name is stored as the value of the
\texttt{svn:author} property on the new revision (see
\hyperref[svn.advanced.props.ref]{Subversion\textquotesingle s Reserved
Properties}). If the client was not authenticated (i.e., if the server
never issued an authentication challenge), the
revision\textquotesingle s \texttt{svn:author} property is empty.

\subsection{Client Credentials}\label{svn.serverconfig.netmodel.creds}

Many Subversion servers are configured to require authentication.
Sometimes anonymous read operations are allowed, while write operations
must be authenticated. In other cases, reads and writes alike require
authentication. Subversion\textquotesingle s different server options
understand different authentication protocols, but from the
user\textquotesingle s point of view, authentication typically boils
down to usernames and passwords. Subversion clients offer several
different ways to retrieve and store a user\textquotesingle s
authentication credentials, from interactive prompting for usernames and
passwords to encrypted and non-encrypted on-disk data caches.

The security-conscious reader will suspect immediately that there is
reason for concern here. ``Caching passwords on disk?
That\textquotesingle s terrible! You should never do that!''
Don\textquotesingle t worry---it\textquotesingle s not as bad as it
sounds. The following sections discuss the various types of credential
caches that Subversion uses, when it uses them, and how to disable that
functionality in whole or in part.

\subsubsection{Caching
credentials}\label{svn.serverconfig.netmodel.credcache}

Subversion offers a remedy for the annoyance caused when users are
forced to type their usernames and passwords over and over again. By
default, whenever the command-line client successfully responds to a
server\textquotesingle s authentication challenge, credentials are
cached on disk and keyed on a combination of the
server\textquotesingle s hostname, port, and authentication realm. This
cache will then be automatically consulted in the future, avoiding the
need for the user to re-type his or her authentication credentials. If
seemingly suitable credentials are not present in the cache, or if the
cached credentials ultimately fail to authenticate, the client will, by
default, fall back to prompting the user for the necessary information.

The Subversion developers recognize that on-disk caches of
authentication credentials can be a security risk. To offset this,
Subversion works with available mechanisms provided by the operating
system and environment to try to minimize the risk of leaking this
information.

\begin{itemize}
\item
  On Windows, the Subversion client stores passwords in the
  \texttt{\%APPDATA\%/Subversion/auth/} directory. On Windows 2000 and
  later, the standard Windows cryptography services are used to encrypt
  the password on disk. Because the encryption key is managed by Windows
  and is tied to the user\textquotesingle s own login credentials, only
  the user can decrypt the cached password. (Note that if the
  user\textquotesingle s Windows account password is reset by an
  administrator, all of the cached passwords become undecipherable. The
  Subversion client will behave as though they don\textquotesingle t
  exist, prompting for passwords when required.)
\item
  Similarly, on Mac OS X, the Subversion client stores all repository
  passwords in the login keyring (managed by the Keychain service),
  which is protected by the user\textquotesingle s account password.
  User preference settings can impose additional policies, such as
  requiring that the user\textquotesingle s account password be entered
  each time the Subversion password is used.
\item
  For other Unix-like operating systems, no single standard ``keychain''
  service exists. However, the Subversion client knows how to store
  passwords securely using the ``GNOME Keyring'', ``KDE Wallet'', and
  ``GnuPG Agent'' services. Also, before storing unencrypted passwords
  in the \texttt{\textasciitilde{}/.subversion/auth/} caching area, the
  Subversion client will ask the user for permission to do so. Note that
  the \texttt{auth/} caching area is still permission-protected so that
  only the user (owner) can read data from it, not the world at large.
  The operating system\textquotesingle s own file permissions protect
  the passwords from other non-administrative users on the same system,
  provided they have no direct physical access to the storage media of
  the home directory, or backups thereof.
\end{itemize}

Of course, for the truly paranoid, none of these mechanisms meets the
test of perfection. So for those folks willing to sacrifice convenience
for the ultimate in security, Subversion provides various ways of
disabling its credentials caching system altogether.

\subsubsection{Disabling password
caching}\label{svn.tour.initial.disabling-password-caching}

When you perform a Subversion operation that requires you to
authenticate, by default Subversion tries to cache your authentication
credentials on disk in encrypted form. On some systems, Subversion may
be unable to encrypt your authentication data. In those situations,
Subversion will ask whether you want to cache your credentials to disk
in plaintext:

\begin{verbatim}
$ svn checkout https://host.example.com:443/svn/private-repo
-----------------------------------------------------------------------
ATTENTION!  Your password for authentication realm:

   <https://host.example.com:443> Subversion Repository

can only be stored to disk unencrypted!  You are advised to configure
your system so that Subversion can store passwords encrypted, if
possible.  See the documentation for details.

You can avoid future appearances of this warning by setting the value
of the 'store-plaintext-passwords' option to either 'yes' or 'no' in
'/tmp/servers'.
-----------------------------------------------------------------------
Store password unencrypted (yes/no)? 
\end{verbatim}

If you want the convenience of not having to continually reenter your
password for future operations, you can answer \texttt{yes} to this
prompt. If you\textquotesingle re concerned about caching your
Subversion passwords in plaintext and do not want to be asked about it
again and again, you can disable caching of plaintext passwords either
permanently, or on a server-by-server basis.

When considering how to use Subversion\textquotesingle s password
caching system, you\textquotesingle ll want to consult any governing
policies that are in place for your client computer---many companies
have strict rules about the ways that their employees\textquotesingle{}
authentication credentials should be stored.

To permanently disable caching of passwords in plaintext, add the line
\texttt{store-plaintext-passwords\ =\ no} to the \texttt{{[}global{]}}
section in the \texttt{servers} configuration file on the local machine.
To disable plaintext password caching for a particular server, use the
same setting in the appropriate group section in the \texttt{servers}
configuration file. (See \hyperref[svn.advanced.confarea.opts]{Runtime
Configuration Options} in \hyperref[svn.customization]{Customizing Your
Subversion Experience} for details.)

To disable password caching entirely for any single Subversion
command-line operation, pass the \texttt{-\/-no-auth-cache} option to
that command line. To permanently disable caching entirely, add the line
\texttt{store-passwords\ =\ no} to your local machine\textquotesingle s
Subversion configuration file.

\subsubsection{Removing cached
credentials}\label{svn.tour.initial.authn-cache-purge}

Sometimes users will want to remove specific credentials from the disk
cache. To do this, you need to navigate into the \texttt{auth/} area and
manually delete the appropriate cache file. Credentials are cached in
individual files; if you look inside each file, you will see keys and
values. The \texttt{svn:realmstring} key describes the particular server
realm that the file is associated with:

\begin{verbatim}
$ ls ~/.subversion/auth/svn.simple/
5671adf2865e267db74f09ba6f872c28
3893ed123b39500bca8a0b382839198e
5c3c22968347b390f349ff340196ed39

$ cat ~/.subversion/auth/svn.simple/5671adf2865e267db74f09ba6f872c28

K 8
username
V 3
joe
K 8
password
V 4
blah
K 15
svn:realmstring
V 45
<https://svn.domain.com:443> Joe's repository
END
\end{verbatim}

Once you have located the proper cache file, just delete it.

\subsubsection{Command-line
authentication}\label{svn.tour.initial.different-user}

All Subversion command-line operations accept the \texttt{-\/-username}
and \texttt{-\/-password} options, which allow you to specify your
username and password, respectively, so that Subversion
isn\textquotesingle t forced to prompt you for that information. This is
especially handy if you need to invoke Subversion from a script and
cannot rely on Subversion being able to locate valid cached credentials
for you. These options are also helpful when Subversion has already
cached authentication credentials for you, but you know they
aren\textquotesingle t the ones you want it to use. Perhaps several
system users share a login to the system, but each have distinct
Subversion identities. You can omit the \texttt{-\/-password} option
from this pair if you wish Subversion to use only the provided username,
but still prompt you for that username\textquotesingle s password.

\subsubsection{Authentication
wrap-up}\label{svn.tour.initial.authn-wrapup}

One last word about \texttt{svn}\textquotesingle s authentication
behavior, specifically regarding the \texttt{-\/-username} and
\texttt{-\/-password} options. Many client subcommands accept these
options, but it is important to understand that using these options does
\emph{not} automatically send credentials to the server. As discussed
earlier, the server ``pulls'' credentials from the client when it deems
necessary; the client cannot ``push'' them at will. If a username and/or
password are passed as options, they will be presented to the server
only if the server requests them. These options are typically used to
authenticate as a different user than Subversion would have chosen by
default (such as your system login name) or when trying to avoid
interactive prompting (such as when calling \texttt{svn} from a script).

A common mistake is to misconfigure a server so that it never issues an
authentication challenge. When users pass \texttt{-\/-username} and
\texttt{-\/-password} options to the client, they\textquotesingle re
surprised to see that they\textquotesingle re never used; that is, new
revisions still appear to have been committed anonymously!

Here is a final summary that describes how a Subversion client behaves
when it receives an authentication challenge.

\begin{enumerate}
\def\labelenumi{\arabic{enumi}.}
\item
  First, the client checks whether the user specified any credentials as
  command-line options (\texttt{-\/-username} and/or
  \texttt{-\/-password}). If so, the client will try to use those
  credentials to authenticate against the server.
\item
  If no command-line credentials were provided, or the provided ones
  were invalid, the client looks up the server\textquotesingle s
  hostname, port, and realm in the runtime
  configuration\textquotesingle s \texttt{auth/} area, to see whether
  appropriate credentials are cached there. If so, it attempts to use
  those credentials to authenticate.
\item
  Finally, if the previous mechanisms failed to successfully
  authenticate the user against the server, the client resorts to
  interactively prompting the user for valid credentials (unless
  instructed not to do so via the \texttt{-\/-non-interactive} option or
  its client-specific equivalents).
\end{enumerate}

If the client successfully authenticates by any of these methods, it
will attempt to cache the credentials on disk (unless the user has
disabled this behavior, as mentioned earlier).

\section{Working Without a Working
Copy}\label{svn.advanced.working-without-a-wc}

As we described in \hyperref[svn.basic.in-action.wc]{Subversion Working
Copies}, the Subversion working copy is a sort of staging area where a
user can privately make changes to his or her versioned data and
then---when those changes are complete and ready for sharing with
others---commit them to the repository. It should come as no surprise,
then, that most of the interaction you will have with Subversion will be
in the form of asking your Subversion client to do \emph{something} to
one or more items in a local working copy. Even for those operations
which don\textquotesingle t manipulate the working copy data itself
(such as \texttt{svn\ log}), it\textquotesingle s often just easier to
use a working copy file or directory as a convenient target for that
operation.

Clearly, the typical approach to making changes to your versioned data
is via commits from a Subversion working copy. Fortunately,
it\textquotesingle s not the only way. Users of Subversion who need to
make relatively simple changes to their versioned data can do so without
the overhead of checking out a working copy. We\textquotesingle ll cover
some of those supported operations in this section.

\subsection{Remote command-line client
operations}\label{svn.advanced.working-without-a-wc.svn}

The Subversion command-line client supports a number of operations which
can be performed directly against repository URLs in order to make
simple changes without a working copy. Some of these are described
elsewhere in this book, but we provide an exhaustive list of them here
for your convenience.

Perhaps the most obvious remote commit-like operation is the
\texttt{svn\ import} command. We describe that command in
\hyperref[svn.tour.importing.import]{Importing Files and Directories} as
part of explaining how you can easily get a whole tree of unversioned
information into your Subversion repository so you can start doing
version-controlled operations on it.

The \texttt{svn\ mkdir} and \texttt{svn\ delete} commands, when used
with URL targets, are also remote commit-type operations. These allow
the user to create one or more new versioned directories or remove
(recursively) one or more versioned files or directories, respectively,
without the use of a working copy. Each time you issue one of these
commands, the client communicates with the server in a way
that\textquotesingle s similar to how it would describe the commit of a
directory added or of an item removed from the working copy. If
there\textquotesingle s no problem or conflict detected with the
requested operation, the server commits the additions or removals in a
single new revision.

You can use \texttt{svn\ copy} or \texttt{svn\ move} with two URLs---a
copy/move source and a destination---to commit a copies and moves of
files and directories directly in the repository. These operations tend
to be some of the most expensive ones when performed within a working
copy, but they complete in constant time when performed remotely using
repository URLs. In fact, the \texttt{svn\ copy} remote operation is
commonly used to create branches in Subversion, as we discuss later in
\hyperref[svn.branchmerge.using.create]{Creating a Branch}.

As with the regular \texttt{svn\ commit} command, you can supply a log
message with any of these commands we\textquotesingle ve discussed so
far to describe the changes you\textquotesingle re making. Use the
\texttt{-\/-file\ (-F)} or \texttt{-\/-message\ (-m)} option, or
otherwise allow the client to prompt you for the log message.

Finally, there are a number of operations related to unversioned
revision properties which can be performed directly against the
repository. In fact, revision properties are somewhat unique in this
context, as they aren\textquotesingle t stored in the working copy and
therefore \emph{must} be modified without working copy interaction. See
\hyperref[svn.advanced.props]{Properties} for a more detailed
description of how to manage properties in Subversion.

\subsection{Using
svnmucc}\label{svn.advanced.working-without-a-wc.svnmucc}

One shortcoming of the remote commit operation support offered in the
command-line client is that you are essentially limited to one
operation---or, really, one type of operation---per commit. For example,
it\textquotesingle s perfectly natural and supported to, say, use
\texttt{svn\ delete} followed by \texttt{svn\ mkdir} within a working
copy to replace an existing versioned directory with a brand new one.
When you commit the results of those operations, a single new revision
is created in the repository, and that revision carries the full
replacement of your directory. You can\textquotesingle t really do the
same thing as remote operations using the command-line client while
still preserving the it-happened-in-a-single-revision-ness of the
change---\texttt{svn\ delete\ URL} would create a new revision that
removed the directory; \texttt{svn\ mkdir\ URL} would generate a second
revision for the directory\textquotesingle s re-creation.

Fortunately, Subversion provides a separate tool which exists solely to
allow users to string together a set of remote operations and commit
them as one atomic change. That tool is the \texttt{svnmucc} tool---the
Subversion Multiple URL Command Client:

\begin{verbatim}
$ svnmucc --help
Subversion multiple URL command client
usage: svnmucc ACTION...

  Perform one or more Subversion repository URL-based ACTIONs, committing
  the result as a (single) new revision.

Actions:
  cp REV URL1 URL2       : copy URL1@REV to URL2
  mkdir URL              : create new directory URL
  mv URL1 URL2           : move URL1 to URL2
  rm URL                 : delete URL
  put SRC-FILE URL       : add or modify file URL with contents copied from
                           SRC-FILE (use "-" to read from standard input)
  propset NAME VAL URL   : set property NAME on URL to value VAL
  propsetf NAME VAL URL  : set property NAME on URL to value from file VAL
  propdel NAME URL       : delete property NAME from URL
…
\end{verbatim}

\texttt{svnmucc} has been a part of the Subversion
project\textquotesingle s source code tree for many years (as
\texttt{mucc} for most of that time), but it was only in Subversion 1.8
that it become a fully supported member of the Subversion command-line
tool suite.

The \texttt{svnmucc} tool can perform any transformation on your
versioned data that \texttt{svn} itself can. But unlike \texttt{svn},
the functionality that \texttt{svnmucc} offers isn\textquotesingle t
broken up into subcommands. Rather, you provide a list of actions and
operands in a single command line (or from a file stream, via the
\texttt{-\/-extra-args\ (-X)} option). Some of the actions supported by
\texttt{svnmucc} mimic those of the command-line client.
You\textquotesingle ll notice in the previous command output actions
such as \texttt{cp}, \texttt{mkdir}, \texttt{mv}, and \texttt{rm}, all
of which are very similar to the commands we mentioned in
\hyperref[svn.advanced.working-without-a-wc.svn]{Remote command-line
client operations}. But remember, the key difference here is that you
can use any number of these actions together in a single command
invocation, resulting in a single committed revision in the repository.

Let\textquotesingle s take our previous example of trying to simply
replace a remote directory. Using \texttt{svnmucc}, you would accomplish
this as follows:

\begin{verbatim}
$ svnmucc rm http://svn.example.com/projects/sandbox \
          mkdir http://svn.example.com/projects/sandbox \
          -m "Replace my old sandbox with a fresh new one."
r22 committed by harry at 2013-01-15T21:45:26.442865Z
$ 
\end{verbatim}

As you can see, \texttt{svnmucc} accomplished in a single revision what
\texttt{svn}---without the benefit of a working copy---required two
revisions to complete.

Another difference between \texttt{svnmucc} and \texttt{svn} is that the
former currently will not prompt you for a commit log message if you
fail to supply one via the command line. Rather, it will use a stock
(that is, relatively valueless) log message.

The \texttt{svnmucc} tool is not limited to merely remixing actions that
\texttt{svn} itself can perform. It introduces some additional
functionality not found in the command-line client. For example, you can
use the \texttt{put} action to add or modify a file in the repository,
copying the file\textquotesingle s intended new contents from either a
file on your local machine or from data piped in via standard input. The
tool also offers \texttt{propset}, \texttt{propsetf}, and
\texttt{propdel} actions, useful for setting properties on versioned
files and directories (explicitly, or by copying the
property\textquotesingle s value from a local file) and for deleting
properties on the same. Those actions are unsupported in the
command-line client at this time.

At this point, though, it seems prudent to discuss the difference
between what \emph{can} be done with \texttt{svnmucc} and what
\emph{should} be done. A pair of notable quotes comes to mind:

\begin{quote}
``To whom much has been given, much will be expected.''

--- Jesus
\end{quote}

\begin{quote}
``With great power comes great responsibility.''

--- "Spiderman" Peter Parker\textquotesingle s Uncle Ben
\end{quote}

Inherent in modifications without a working copy is the loss of the very
conflict detection safeguards which make the use of a working copy so
valuable. When using \texttt{svn} in the typical way, changes are
committed to the server against a specific base version of a file or
directory so that you don\textquotesingle t inadvertently overwrite
contemporary changes made to the same item by another team member. The
server knows what version of the file you had before you changed it, and
it knows if other folks have changed that same file since that revision
was created. That\textquotesingle s all the information the server needs
to deny your commit when it would clobber someone else\textquotesingle s
change, forcing you to integrate their change into your working copy and
reconsider your own change. Because there is no working copy in the mix
here, \texttt{svnmucc} really gives you the power to bypass those
safeguards and to act as if the current state of the repository is
precisely the base state against which you are working. But hopefully it
is obvious to you that this is not a power you should cavalierly wield.

Fortunately, \texttt{svnmucc} allows you to be more conservative in the
way you use the tool. In order to provide a safety mechanism similar to
what is offered by the use of a working copy, \texttt{svnmucc} offers a
\texttt{-\/-revision\ (-r)} option. With this option, you can manually
specify a base revision for the changes you are attempting to commit.
The base revision you choose is ideally the most recent revision in your
repository of which you can reasonably claim knowledge.

Users are strongly encouraged to use, and to use correctly, the
\texttt{-\/-revision\ (-r)} option to \texttt{svnmucc}.

Proper use of the \texttt{svnmucc\ put} action best demonstrates how
this \texttt{-\/-revision\ (-r)} option should be used. Say Harry wishes
to change the contents of a versioned \texttt{README} file without
bothering with a full checkout of a working copy. (We\textquotesingle ll
assume that there is no other value in using a working copy for this
operation, such as the presence of scripts Harry should run in advance
of his commit to verify that it\textquotesingle s a reasonable one.) The
first decision he has to make is which revision of the file he wants to
work with. Typically, users wish to modify the most recent version of a
file. So Harry queries the revision in which the file was last modified,
and then uses that revision to fetch the contents of the file into a
temporary local file:

\begin{verbatim}
$ svn info http://svn.example.com/projects/sandbox/README
Path: README
URL: http://svn.example.com/projects/sandbox/README
Relative URL: ^/sandbox/README
Repository Root: http://svn.example.com/projects
Repository UUID: 13f79535-47bb-0310-9956-ffa450edef68
Revision: 22
Node Kind: file
Last Changed Author: sally
Last Changed Rev: 14
Last Changed Date: 2012-09-02 10:34:09 -0400 (Sun, 02 Sep 2012)

$ svn cat -r 14 http://svn.example.com/projects/sandbox/README \
      > README.tmpfile
$
\end{verbatim}

Harry now has a copy of the \texttt{README} file as it looked when it it
was last modified. He makes the edits he wishes to make to this copy of
the file. Naturally, when he\textquotesingle s finished, he wishes to
then commit those changes to the repository.

Now, if Harry naively uses \texttt{svnmucc\ put\ …} at this point to
replace the contents of \texttt{README} in the repository with his
locally modified contents, he has just abused the power that
\texttt{svnmucc} affords. What if, just microseconds prior to his
commit, Sally had also modified the \texttt{README} file? As with the
\texttt{svn} program, \texttt{svnmucc} won\textquotesingle t attempt
some sort of server-side content merge in order to preserve both
users\textquotesingle{} changes. Rather, \texttt{svnmucc} will happily
replace the current latest version of the file with the contents
specified. Harry will be oblivious. Sally will be livid.

\begin{verbatim}
$ svnmucc put README.tmpfile \
          http://svn.example.com/projects/sandbox/README \
          -m "Tweak the README file."
r24 committed by harry at 2013-01-21T16:21:23.100133Z
$
Message from sally@shell.example.com on pts/2 at 16:26 ...
We need to talk.  Now.
EOF
\end{verbatim}

Harry should instead recall the revision he originally used as the
revision on which to base his changes, supplying that revision to
\texttt{svnmucc} via the \texttt{-\/-revision\ (-r)} option, and thus
giving the server the opportunity to bounce his commit if, by his own
(perhaps ignorant) admission, he\textquotesingle s attempting to modify
an out-of-date item:

\begin{verbatim}
$ svnmucc -r 14 put README.tmpfile \
          http://svn.example.com/projects/sandbox/README \
          -m "Tweak the README file."
svnmucc: E170004: Item '/sandbox/README' is out of date
$
\end{verbatim}

Like other \texttt{svnmucc} options, the \texttt{-\/-revision\ (-r)}
option operates at a scope global to the whole command---every action
specified in that command. This enables you to have the same sort of
safeguards you would have if you had checked out a working copy of your
entire repository (and thus had a working copy entirely at a single
uniform revision), made changes to that working copy, and then committed
all those changes at once.

As you can see, \texttt{svnmucc} is a handy addition to the Subversion
user\textquotesingle s tool chest. For a complete reference of this
tool\textquotesingle s offerings, see \hyperref[svn.ref.svnmucc]{svnmucc
Reference---Subversion Multiple URL Command Client}.

\section{Summary}\label{svn.advanced.summary}

After reading this chapter, you should have a firm grasp on some of
Subversion\textquotesingle s features that, while perhaps not used
\emph{every} time you interact with your version control system, are
certainly handy to know about. But don\textquotesingle t stop here! Read
on to the following chapter, where you\textquotesingle ll learn about
branches, tags, and merging. Then you\textquotesingle ll have nearly
full mastery of the Subversion client. Though our lawyers
won\textquotesingle t allow us to promise you anything, this additional
knowledge could make you measurably more cool.\footnote{No purchase
  necessary. Certains terms and conditions apply. No guarantee of
  coolness---implicit or otherwise---exists. Mileage may vary.}

